%==============================================================================
% Chapter 1 — Introduction
%==============================================================================
\chapter{المقدمة}
\label{ch:introduction}
\markboth{الفصل الأول: المقدمة}{الفصل الأول: المقدمة}

%------------------------------------------------------------------------------
\section{الخلفية (\EN{Background})}
\label{sec:background}

شهد العالم خلال السنوات الأخيرة تحولًا كبيرًا في طريقة إدارة المؤسسات، حيث
أصبح الاعتماد على الأنظمة الرقمية والذكاء الاصطناعي عاملًا رئيسيًا في تحسين
جودة الخدمات المقدمة، وتعزيز كفاءة التشغيل، وتقليل الجهد اليدوي. وفقاً
لتقرير \EN{McKinsey Global Institute} لعام 2024 \cite{mckinsey2024ai}، من المتوقع أن يُسهم الذكاء
الاصطناعي بخلق قيمة اقتصادية تتراوح بين 17 و25 تريليون دولار سنوياً بحلول
عام 2030، يُخصص منها نحو 40\% لتطبيقات المؤسسات والإدارة. ومع النمو المتسارع
في حجم البيانات والمستندات المؤسسية — الذي يُقدر بـ 175 زيتابايت عالمياً
بحلول 2025 وفق توقعات \EN{IDC} \cite{idc2025datasphere} — وازدياد عدد المستخدمين والخدمات المرتبطة
بها، أصبحت إدارة المعلومات أكثر تعقيدًا من أي وقت مضى، ولم تعد الأساليب
التقليدية قادرة على مواكبة حجم العمل ومتطلبات المستخدمين.

في المؤسسات الحديثة، تتنوع المهام اليومية التي تتولى الإدارة الإشراف عليها،
بدءًا من استعلامات البيانات وتحليلها، مرورًا بالبحث في المستندات واللوائح،
ووصولاً إلى إعداد التقارير واتخاذ القرارات الاستراتيجية. هذه المهام تحتاج
إلى نظام مركزي موثوق يضمن تنظيمها وتوثيقها، ويمنع ضياع المعلومات أو
تداخلها، خاصة مع اعتماد المستخدمين المتزايد على الخدمات السريعة والموثوقة.

ورغم هذا التطور العالمي، ما زالت الكثير من المؤسسات في اليمن وفي العديد من
الدول العربية تُدار بطرق تقليدية تُعرقل سير العمل وتؤثر على كفاءة اتخاذ
القرار. ما يزال الاعتماد على الاستعلامات المباشرة لقواعد البيانات من قبل
متخصصين شائعًا، وكذلك البحث اليدوي في المستندات غير المفهرسة، الأمر الذي
يفتقر إلى السرعة والدقة ويؤدي إلى ضياع الوقت أو تكرار الأخطاء. كما أن إعداد
التقارير والتحليلات تقتصر غالبًا على الجهد اليدوي دون وجود أي أدوات تحليلية
تمكّن المستخدم من فهم اتجاهات البيانات أو اتخاذ قرارات مبنية على المعلومات.

ومع غياب نظام موحّد يجمع كل هذه الوظائف، تصبح عملية إدارة المعلومات مرهقة
وغير فعّالة، ويعاني المستخدم من نقص الشفافية وضعف التواصل مع إدارة المؤسسة.
كما تواجه الإدارة تحديات مستمرة في متابعة استعلامات المستخدمين، والتحقق من
دقة المعلومات، وتحديد احتياجات التدريب، وإصدار تقارير تساعد في تحسين مستوى
الخدمة.

من جهة أخرى، أدى التقدم في مجال الذكاء الاصطناعي وتعلم الآلة إلى فتح الباب
أمام حلول رقمية جديدة في مجال إدارة المؤسسات، سواء عبر المساعدات الذكية أو
أنظمة المحادثة المتقدمة، مما جعل عملية التحول الرقمي أكثر قابلية للتطبيق حتى
في البيئات محدودة الموارد. هذه الحلول أصبحت اليوم ضرورة وليس خيارًا، خاصة في
ظل احتياج المستخدمين إلى خدمات أسرع وأكثر دقة وأعلى جودة.

وبناءً على ذلك، نشأت الحاجة إلى تطوير نظام شامل لإدارة المؤسسات يربط بين
المستخدمين والإدارة في إطار واحد، ويتيح تبادل المعلومات بشكل فوري، ويعزز
الشفافية، ويوفّر سجلًا موحدًا لكل العمليات اليومية. ومن هنا جاء مشروع
«المساعد الذكي المؤسسي (بيان)» ليقدّم نظامًا ذكيًا مصممًا خصيصًا لمعالجة
التحديات التي تواجه المؤسسات في الواقع المحلي، مع مراعاة سهولة الاستخدام،
وقابلية التوسع، والقدرة على تقديم تجربة استخدام احترافية تجمع بين البساطة
والفعالية.

يقدم نظام «بيان» خطوة مهمة نحو التحول الرقمي في إدارة المؤسسات، حيث يجمع
بين أحدث تقنيات الذكاء الاصطناعي وتطوير التطبيقات وبين فهم عميق لاحتياجات
المستخدمين والإدارة، مما يجعله حلًا مبتكرًا وعمليًا يمكن تطبيقه مباشرة في
المؤسسات، ويسهم في رفع مستوى الخدمات وتحسين كفاءة العمل واتخاذ القرارات.

%------------------------------------------------------------------------------
\section{مشكلة المشروع (\EN{Problem Statement})}
\label{sec:problem}

تكمن المشكلة الرئيسية في تعثر الوصول الذكي والسريع إلى البيانات والمستندات
المؤسسية نتيجة الاعتماد على الوسائل التقليدية واليدوية، مما يترتب عليه بطء
في اتخاذ القرار، وهدر الموارد التشغيلية، وانخفاض مستوى جودة الخدمات المقدمة.

\subsection*{مظاهر المشكلة}

\begin{enumerate}
  \item \textbf{صعوبة الوصول المباشر للبيانات:} اقتصار استعلام قواعد البيانات
        على المتخصصين التقنيين (\EN{SQL})، مما يحرم الإدارة والموظفين غير
        المختصين من الاستعلام المباشر ويؤخر الحصول على المعلومة.
  \item \textbf{شتات المستندات المؤسسية:} غياب آلية موحدة للبحث في المستندات
        والسياسات واللوائح غير المفهرسة، مما يجعل الوصول إلى نصوص القرارات
        واللوائح أمراً مستغرقاً للوقت.
  \item \textbf{بطء وتكلفة إعداد التقارير:} الاعتماد على تجميع البيانات
        وإعداد التقارير والتحليلات يدوياً، مما يزيد من احتمالية الأخطاء
        التشغيلية ويؤخر صُنّاع القرار.
  \item \textbf{غياب الشفافية وسجلات التتبع:} تفتقر العديد من العمليات
        التشغيلية إلى سجل تتبع موحد للاستعلامات والإجابات، مما يعوق تقييم
        الأداء والرقابة.
  \item \textbf{ضعف التواصل المؤسسي:} غياب القنوات الذكية المباشرة للتفاعل
        بين المستفيدين والإدارة، مما ينعكس سلباً على الشفافية ورضا المستخدمين.
\end{enumerate}

%------------------------------------------------------------------------------
\section{أهداف المشروع (\EN{Objectives})}
\label{sec:objectives}

يهدف مشروع «بيان» إلى تطوير نظام مؤسسي متكامل مدعوم بتقنيات الذكاء الاصطناعي
وتحليل البيانات، لتمكين المؤسسات من الوصول الذكي إلى بياناتها ومستنداتها
باللغة العربية الطبيعية، وتحويل الأسئلة والمستندات إلى إجابات وتقارير دقيقة
وفورية. ويتفرع من هذا الهدف الرئيسي مجموعة من الأهداف التفصيلية الآتية:

\begin{itemize}
  \item \textbf{بناء مساعد محادثة ذكي باللغة العربية:} التفاعل مع المستخدمين
        باللغة العربية الطبيعية (الفصحى والعامية).
  \item \textbf{تمكين الاستعلام الذكي عن البيانات} (\EN{NL2SQL}): تحويل
        الاستعلامات الصادرة باللغة الطبيعية إلى استعلامات قواعد بيانات
        (\EN{SQL}) دون حاجة المستخدم لمعرفة تقنية مسبقة.
  \item \textbf{تطوير نظام استرجاع المستندات} (\EN{RAG}): البحث في المستندات
        واللوائح المؤسسية وتوليد إجابات دقيقة وموثقة بمصادرها.
  \item \textbf{توفير أدوات التحليل الذكي:} تحليل البيانات واستخراج الاتجاهات
        والقيم الشاذة وتقديم رؤى إحصائية تدعم اتخاذ القرار.
  \item \textbf{توليد التقارير المنسقة:} إتاحة إمكانية إنتاج تقارير متعددة
        الأشكال (نصية، جدولية، ورسوم بيانية) قابلة للتصدير بصيغ مختلفة.
  \item \textbf{ضمان الشفافية والرقابة:} توفير سجل موحد (\EN{Audit Log})
        لتتبع الاستعلامات والإجابات وسجلات الوصول تعزيزًا للمساءلة.
  \item \textbf{تحسين التواصل المؤسسي:} تعزيز التفاعل بين المستخدمين وإدارة
        المؤسسة عبر نظام إشعارات وتنبيهات داخلية موحدة.
\end{itemize}

%------------------------------------------------------------------------------
\section{أهمية المشروع (\EN{Significance})}
\label{sec:significance}

تنبع أهمية مشروع «بيان» من كونه يستهدف سد فجوة حقيقية في البيئة المؤسسية
العربية عامةً، واليمنية خاصةً، تتمثل في غياب الأنظمة الذكية الموحّدة التي تجمع
بين الاستعلام عن البيانات والبحث في المستندات وتحليلها في إطار واحد. فالأنظمة
التجارية المتاحة تُعرَض بأسعار مرتفعة وتتطلب بنية تحتية متقدمة، كما أنها لا
تدعم اللغة العربية بالشكل المطلوب، مما يُقصي غالبية المؤسسات المحلية عن
الاستفادة من هذه التقنيات.

على الصعيد التشغيلي، يُتوقع أن يسهم النظام في تسريع عملية اتخاذ القرار من
خلال تقصير زمن الاستعلام عن المعلومة من دقائق أو ساعات (في الوضع اليدوي
التقليدي) إلى ثوانٍ معدودة. كما يُتوقع أن يُقلّل من الأخطاء البشرية الناتجة عن
التعامل اليدوي مع البيانات، ويرفع من مستوى الشفافية عبر سجل التدقيق الشامل
(\EN{Audit Log}) الذي يُوثّق كل استعلام وإجابة.

على الصعيد الأكاديمي، يُمثّل المشروع إضافة معرفية في مجال دمج التقنيات الأربع
(\EN{NL2SQL}، \EN{RAG}، \EN{Multi-Agent}، \EN{Arabic NLP}) في إطار موحد، وهو
حقل بحثي لم يُطرق بكثرة في الأدبيات العربية. كما يفتح الباب أمام دراسات لاحقة
لتطوير النظام ليشمل اللهجات المحلية، أو توسيع نطاقه ليشمل مؤسسات حكومية أخرى.

على الصعيد المجتمعي، يُسهم المشروع في دعم التحول الرقمي الذي تتبناه الدولة
في إطار رؤية التحول الرقمي اليمني \cite{un2024egd}، ويعزز من قدرة المؤسسات على تقديم خدمات
أفضل للمواطنين عبر توفير معلومات دقيقة وسريعة، مما ينعكس إيجاباً على رضا
المستفيدين وثقتهم في الخدمات الحكومية.

%------------------------------------------------------------------------------
\section{مساهمات المشروع (\EN{Project Contributions})}
\label{sec:contributions}

يُقدّم مشروع «بيان» مساهمات بحثية وتطبيقية في أربعة محاور رئيسية، يمكن
البناء عليها في دراسات وأعمال مستقبلية:

\begin{enumerate}
  \item \textbf{المساهمة التقنية:} دمجُ أربع تقنيات متقدمة (\EN{NL2SQL}،
        \EN{RAG}، \EN{Multi-Agent}، \EN{Arabic NLP}) في إطار معماري موحد عبر
        \EN{LangGraph}، وهو تكامل غير مسبوق في الأدبيات العربية وغير معتاد حتى
        في الأدبيات الإنجليزية، حيث تُقدَّم هذه التقنيات عادةً بشكل منفصل.
  \item \textbf{المساهمة الابتكارية:} ابتكار طبقة الرؤى الذكية (\EN{Smart Views}) كآلية
        لتبسيط مخطط قاعدة البيانات قبل توليد \EN{SQL}، مما يرفع دقة
        \EN{NL2SQL} بشكل ملحوظ ويُمثّل مساهمة بحثية أصلية قابلة للتعميم على
        أي نظام إدارة قواعد بيانات مؤسسية.
  \item \textbf{المساهمة اللغوية:} تقديم أول نظام \EN{NL2SQL}+\EN{RAG} عربي
        مفتوح المصدر يدعم الفصحى والعامية في واجهة محادثة \EN{RTL}، مما يسدّ
        فجوةً واضحةً في دعم اللغة العربية في أنظمة الذكاء الاصطناعي المؤسسية.
  \item \textbf{المساهمة التطبيقية:} تقديم تطبيق ميداني متكامل للبنية متعددة
        الوكلاء في بيئة مؤسسية تعليمية، باستخدام نظام \EN{Odoo Education}
        كبيئة تطبيق واقعية، مع توثيق كامل للمنهجية وتقييم الأداء، مما يوفر
        مرجعاً تطبيقياً للمؤسسات العربية الراغبة في تبني تقنيات مشابهة.
\end{enumerate}

%------------------------------------------------------------------------------
\section{تعريف المشروع (\EN{System Definition})}
\label{sec:definition}

يُعرَّف نظام «المساعد الذكي المؤسسي (بيان)» بأنه منصة محادثة ذكية متقدمة قائمة
على معمارية الأنظمة متعددة الوكلاء (\EN{Multi-Agent System})، تمكّن المستخدمين
من التفاعل التكاملي مع قواعد البيانات والمستندات المؤسسية باللغة العربية
الطبيعية. وتتكون البنية الأساسية للنظام من المكونات التالية:

\begin{itemize}
  \item \textbf{واجهة محادثة عربية} (\EN{RTL}): واجهة تفاعلية تدعم اتجاه
        الكتابة من اليمين إلى اليسار لطرح الأسئلة واستقبال الردود الفورية.
  \item \textbf{لوحة تحكم ويب} (\EN{Web Admin Dashboard}): منصة إدارية لمراقبة
        أداء النظام، وتتبع الوكلاء، وإدارة الصلاحيات والأمان.
  \item \textbf{خمسة وكلاء ذكاء اصطناعي متخصصين} يعملون في إطار منسق عبر
        تقنية \EN{LangGraph}، وهم:
    \begin{itemize}
      \item الوكيل المشرف (\EN{Supervisor Agent}): لتوجيه الاستعلامات
            وإدارة تدفق العمل بين الوكلاء.
      \item وكيل \EN{SQL} (\EN{SQL Agent}): لتحويل اللغة الطبيعية إلى
            استعلامات قواعد بيانات.
      \item وكيل \EN{RAG} (\EN{RAG Agent}): للاسترجاع المعزز بالتوليد
            والبحث في المستندات.
      \item وكيل التحليل (\EN{Analyst Agent}): لاستخراج الرؤى والاتجاهات
            الإحصائية من البيانات.
      \item وكيل التقارير (\EN{Reporter Agent}): لتنسيق المخرجات وتوليد
            التقارير بصيغ متعددة.
    \end{itemize}
  \item \textbf{قاعدة بيانات موحدة وآمنة:} لحفظ سجلات المستخدمين، والأنشطة،
        والمحادثات المخزنة.
  \item \textbf{طبقة الرؤى الذكية} (\EN{Smart Views}): طبقة وسيطة مخصصة
        تبسط بنية قاعدة البيانات لرفع دقة تحويل النص إلى استعلامات
        (\EN{NL2SQL}).
  \item \textbf{بيئة محاكاة \EN{Odoo ERP}:} نظام \EN{ERP} مفتوح المصدر يُستخدم
        كبيئة تطبيق واقعية، يُحقن ببيانات اصطناعية تعادل خمس سنوات من النشاط المؤسسي
        لاختبار النظام في ظروف قريبة من الواقع.
\end{itemize}

ويعمل النظام على رقمنة عمليات الوصول للبيانات وأتمتتها من خلال الآليات
التالية:

\begin{enumerate}
  \item \textbf{المحادثة التفاعلية الذكية:} معالجة اللغة الطبيعية العربية وفهم
        سياق الأسئلة.
  \item \textbf{تحويل النص إلى استعلام} (\EN{NL2SQL}): ترجمة أسئلة المستخدمين
        مباشرة إلى استعلامات قواعد بيانات تنفيذية.
  \item \textbf{الاسترجاع المعزز بالتوليد} (\EN{RAG}): الفهرسة والبحث النصي
        الذكي في المستندات الرسمية وإسناد الإجابات للمصادر.
  \item \textbf{التحليل الإحصائي الآلي:} استخراج المؤشرات، والاتجاهات العامة،
        والانحرافات الإحصائية من البيانات.
  \item \textbf{توليد التقارير:} صياغة وتصدير تقارير الجداول والرسوم البيانية
        بصيغ متعددة.
  \item \textbf{التهيئة الذكية} (\EN{Smart Onboarding}): عملية مؤتمتة تربط النظام
        ببيانات المؤسسة وفهرسة وثائقها خلال زمن قدره 8 ساعات.
  \item \textbf{الأمان المتدرج:} تصنيف أمن الجداول إلى أربعة مستويات أمنية
        (\EN{Public}، \EN{Masked}، \EN{Restricted}، \EN{Blacklist}).
  \item \textbf{تتبع مسار التنفيذ} (\EN{Agent Trace}): الشفافية الكاملة في
        إظهار الخطوات التحليلية التي اتبعها الوكلاء للوصول إلى النتيجة.
\end{enumerate}

%------------------------------------------------------------------------------
\section{نطاق المشروع (\EN{Project Scope})}
\label{sec:scope}

يحدد هذا القسم النطاق الجغرافي والوظيفي للمشروع، بالإضافة إلى أصحاب المصلحة
الرئيسيين المستفيدين من النظام.

\subsection{النطاق الجغرافي (\EN{Geographical Scope})}
يستهدف المشروع المؤسسات والجهات الأكاديمية بشكل عام. وسيُطبق النظام في
مرحلته الأولى على نظام \EN{Odoo Education} المخصص لإدارة البيانات الجامعية،
كبيئة تطبيق واقعية، مع إمكانية تعميم النظام لاحقاً على أي مؤسسة عربية تطلب
بيئة محادثة ذكية باللغة العربية للوصول إلى بياناتها ومستنداتها.

\subsection{النطاق الوظيفي (\EN{Functional Scope})}
يشمل النطاق الوظيفي تطوير نظام محادثة ذكي يدمج بين أربع تقنيات أساسية:
تحويل اللغة الطبيعية إلى \EN{SQL}، والاسترجاع المعزز بالتوليد، والأنظمة
متعددة الوكلاء، ومعالجة اللغة العربية. ويقدم النظام للأنظمة الجامعية حلاً
شاملاً للاستعلام عن البيانات والبحث في المستندات وتوليد التقارير، مع ضمان
الأمان والشفافية.

\subsection{أصحاب المصلحة (\EN{Stakeholders})}
تشمل قائمة أصحاب المصلحة الرئيسيين:
\begin{itemize}
  \item \textbf{صانعو القرار المؤسسي (المديرون التنفيذيون):} يحتاجون رؤى
        استراتيجية سريعة وموثقة.
  \item \textbf{محللو البيانات:} يجرون استعلامات معقدة على البيانات.
  \item \textbf{موظفو العمليات:} يبحثون عن معلومات محددة بسرعة.
  \item \textbf{مسؤولو النظام:} يشرفون على التشغيل والأمان.
  \item \textbf{مشغلو التهيئة:} مسؤولون عن الإعداد الأولي للمؤسسة.
  \item \textbf{الإدارة الأكاديمية:} تتبني النظام كأداة تحول رقمي.
\end{itemize}

%------------------------------------------------------------------------------
\section{قيود المشروع (\EN{Limitations})}
\label{sec:limitations}

يواجه المشروع مجموعة من القيود التي يجب أخذها بعين الاعتبار:
\begin{itemize}
  \item \textbf{قيود البنية التحتية:} محدودية الموارد الحاسوبية محلياً، مما
        يستوجب الاعتماد على نماذج لغوية متوسطة الحجم.
  \item \textbf{قيود اللغة:} ضعف دعم اللغة العربية في النماذج اللغوية
        التجارية، مما يستدعي حلولاً تكميلية.
  \item \textbf{قيود البيانات:} ندرة قواعد البيانات العربية الكاملة للاختبار.
  \item \textbf{قيود الوقت:} كون المشروع ضمن متطلبات التخرج يفرض جدولاً زمنياً
        محدوداً.
  \item \textbf{قيود التمويل:} عدم توفر ميزانية للاعتماد على نماذج سحابية
        مدفوعة، مما يُلزم بالاعتماد على حلول مفتوحة المصدر.
  \item \textbf{قيود الكوادر:} محدودية الفريق المعني بالتطوير والصيانة.
\end{itemize}

%------------------------------------------------------------------------------
\section{المنهجية المتبعة في تطوير المشروع (\EN{Agile Methodology})}
\label{sec:agile}

\begin{figure}[H]
  \centering
  \resizebox{0.7\textwidth}{!}{%==============================================================================
% figures/tikz/agile-cycle.tex
% Professional Agile methodology cycle diagram for Bayan project.
% Renders natively in LaTeX (real text, Amiri font, vector PDF output).
% Uses the document's color palette (kauprimary, kausecondary, kauaccent,
% kaulight, kaudark) so it visually matches the rest of the thesis.
%
% Usage in any chapter:
%   \begin{figure}[!htbp]
%     \centering
%     \resizebox{0.7\textwidth}{!}{%==============================================================================
% figures/tikz/agile-cycle.tex
% Professional Agile methodology cycle diagram for Bayan project.
% Renders natively in LaTeX (real text, Amiri font, vector PDF output).
% Uses the document's color palette (kauprimary, kausecondary, kauaccent,
% kaulight, kaudark) so it visually matches the rest of the thesis.
%
% Usage in any chapter:
%   \begin{figure}[!htbp]
%     \centering
%     \resizebox{0.7\textwidth}{!}{%==============================================================================
% figures/tikz/agile-cycle.tex
% Professional Agile methodology cycle diagram for Bayan project.
% Renders natively in LaTeX (real text, Amiri font, vector PDF output).
% Uses the document's color palette (kauprimary, kausecondary, kauaccent,
% kaulight, kaudark) so it visually matches the rest of the thesis.
%
% Usage in any chapter:
%   \begin{figure}[!htbp]
%     \centering
%     \resizebox{0.7\textwidth}{!}{\input{figures/tikz/agile-cycle.tex}}
%     \caption{...}
%     \label{...}
%   \end{figure}
%==============================================================================

% Radii (in cm):
%   Ri = inner radius (where chevron arrows start)
%   Ro = outer radius (where chevron arrows end)
%   Tip = chevron tip extension beyond Ro
%   RL  = label radius (outer ring text)
\def\RingRi{2.4}
\def\RingRo{3.6}
\def\RingTip{0.35}
\def\RingRL{4.95}

\begin{tikzpicture}[
  % Six stage color styles — drawn from the Bayan palette family.
  c1/.style={fill=kauprimary!85,   draw=kauprimary,   text=white},
  c2/.style={fill=kausecondary!85, draw=kausecondary, text=white},
  c3/.style={fill=kaudark!85,      draw=kaudark,      text=white},
  c4/.style={fill=kauaccent!85,    draw=kauaccent,    text=white},
  c5/.style={fill=kauprimary!55,   draw=kauprimary,   text=white},
  c6/.style={fill=kausecondary!50, draw=kausecondary, text=white},
  badge/.style={
    circle, fill=white, draw=kauprimary, line width=1pt,
    text=kauprimary, font=\large\bfseries, inner sep=0pt,
    minimum size=1.1cm,
  },
  stage label/.style={
    font=\large\bfseries, text=kaudark, align=center,
  },
  center title/.style={
    font=\Huge\bfseries\color{kauprimary}, align=center,
  },
  center subtitle/.style={
    font=\normalsize\color{kausecondary}, align=center,
  },
  >=stealth,
]

  %------------------------------------------------------------------------
  % Six chevron-arrow segments arranged in a clockwise ring.
  % Each segment spans 60 degrees. Math convention: 0°=right, 90°=top.
  % Stage 1: 30°→90° (top-right)        — Planning
  % Stage 2: -30°→30° (right)            — Design
  % Stage 3: -90°→-30° (bottom-right)    — Development
  % Stage 4: -150°→-90° (bottom-left)    — Testing
  % Stage 5: 150°→210° (left)            — Deployment  [210° = -150°]
  % Stage 6: 90°→150° (top-left)         — Review
  %------------------------------------------------------------------------

  % --- Stage 1: Planning (top-right, 30 to 90 deg) ---
  \fill[c1]
    (90:\RingRi) --
    (90:\RingRo) --
    (35:\RingRo) --
    (30:{\RingRo+\RingTip}) --
    (25:\RingRo) --
    (30:\RingRi) -- cycle;

  % --- Stage 2: Design (right, -30 to 30 deg) ---
  \fill[c2]
    (30:\RingRi) --
    (30:\RingRo) --
    (-25:\RingRo) --
    (-30:{\RingRo+\RingTip}) --
    (-35:\RingRo) --
    (-30:\RingRi) -- cycle;

  % --- Stage 3: Development (bottom-right, -90 to -30 deg) ---
  \fill[c3]
    (-30:\RingRi) --
    (-30:\RingRo) --
    (-85:\RingRo) --
    (-90:{\RingRo+\RingTip}) --
    (-95:\RingRo) --
    (-90:\RingRi) -- cycle;

  % --- Stage 4: Testing (bottom-left, -150 to -90 deg) ---
  \fill[c4]
    (-90:\RingRi) --
    (-90:\RingRo) --
    (-145:\RingRo) --
    (-150:{\RingRo+\RingTip}) --
    (-155:\RingRo) --
    (-150:\RingRi) -- cycle;

  % --- Stage 5: Deployment (left, 150 to 210 deg) ---
  \fill[c5]
    (-150:\RingRi) --
    (-150:\RingRo) --
    (155:\RingRo) --
    (150:{\RingRo+\RingTip}) --
    (145:\RingRo) --
    (150:\RingRi) -- cycle;

  % --- Stage 6: Review (top-left, 90 to 150 deg) ---
  \fill[c6]
    (150:\RingRi) --
    (150:\RingRo) --
    (95:\RingRo) --
    (90:{\RingRo+\RingTip}) --
    (85:\RingRo) --
    (90:\RingRi) -- cycle;

  %------------------------------------------------------------------------
  % Inner decorative rings (subtle outlines for polish).
  %------------------------------------------------------------------------
  \draw[kauprimary!30, line width=0.6pt] (0,0) circle ({\RingRi-0.15});
  \draw[kauprimary!20, line width=0.4pt] (0,0) circle ({\RingRo+\RingTip+0.15});

  %------------------------------------------------------------------------
  % Center title.
  %------------------------------------------------------------------------
  \node[center title] at (0,0.35) {\RL{منهجية}};
  \node[center title] at (0,-0.45) {\RL{أجايل}};
  \node[center subtitle] at (0,-1.15) {\RL{\EN{Agile Methodology}}};

  %------------------------------------------------------------------------
  % Stage labels (Arabic name + English gloss) outside each segment.
  % Position at the mid-angle of each 60-degree segment.
  %------------------------------------------------------------------------
  \node[stage label] at (60:\RingRL)  {\RL{التخطيط}\\[2pt]{\small\color{kausecondary}\EN{Planning}}};
  \node[stage label] at (0:\RingRL)   {\RL{التصميم}\\[2pt]{\small\color{kausecondary}\EN{Design}}};
  \node[stage label] at (-60:\RingRL) {\RL{التطوير}\\[2pt]{\small\color{kausecondary}\EN{Development}}};
  \node[stage label] at (-120:\RingRL){\RL{الاختبار}\\[2pt]{\small\color{kausecondary}\EN{Testing}}};
  \node[stage label] at (180:\RingRL) {\RL{الإطلاق}\\[2pt]{\small\color{kausecondary}\EN{Deployment}}};
  \node[stage label] at (120:\RingRL) {\RL{المراجعة}\\[2pt]{\small\color{kausecondary}\EN{Review}}};

  %------------------------------------------------------------------------
  % Numbered badges (1–6) on each segment, centered on the segment.
  %------------------------------------------------------------------------
  \node[badge] at (60:3.0)  {1};
  \node[badge] at (0:3.0)   {2};
  \node[badge] at (-60:3.0) {3};
  \node[badge] at (-120:3.0){4};
  \node[badge] at (180:3.0) {5};
  \node[badge] at (120:3.0) {6};

\end{tikzpicture}
}
%     \caption{...}
%     \label{...}
%   \end{figure}
%==============================================================================

% Radii (in cm):
%   Ri = inner radius (where chevron arrows start)
%   Ro = outer radius (where chevron arrows end)
%   Tip = chevron tip extension beyond Ro
%   RL  = label radius (outer ring text)
\def\RingRi{2.4}
\def\RingRo{3.6}
\def\RingTip{0.35}
\def\RingRL{4.95}

\begin{tikzpicture}[
  % Six stage color styles — drawn from the Bayan palette family.
  c1/.style={fill=kauprimary!85,   draw=kauprimary,   text=white},
  c2/.style={fill=kausecondary!85, draw=kausecondary, text=white},
  c3/.style={fill=kaudark!85,      draw=kaudark,      text=white},
  c4/.style={fill=kauaccent!85,    draw=kauaccent,    text=white},
  c5/.style={fill=kauprimary!55,   draw=kauprimary,   text=white},
  c6/.style={fill=kausecondary!50, draw=kausecondary, text=white},
  badge/.style={
    circle, fill=white, draw=kauprimary, line width=1pt,
    text=kauprimary, font=\large\bfseries, inner sep=0pt,
    minimum size=1.1cm,
  },
  stage label/.style={
    font=\large\bfseries, text=kaudark, align=center,
  },
  center title/.style={
    font=\Huge\bfseries\color{kauprimary}, align=center,
  },
  center subtitle/.style={
    font=\normalsize\color{kausecondary}, align=center,
  },
  >=stealth,
]

  %------------------------------------------------------------------------
  % Six chevron-arrow segments arranged in a clockwise ring.
  % Each segment spans 60 degrees. Math convention: 0°=right, 90°=top.
  % Stage 1: 30°→90° (top-right)        — Planning
  % Stage 2: -30°→30° (right)            — Design
  % Stage 3: -90°→-30° (bottom-right)    — Development
  % Stage 4: -150°→-90° (bottom-left)    — Testing
  % Stage 5: 150°→210° (left)            — Deployment  [210° = -150°]
  % Stage 6: 90°→150° (top-left)         — Review
  %------------------------------------------------------------------------

  % --- Stage 1: Planning (top-right, 30 to 90 deg) ---
  \fill[c1]
    (90:\RingRi) --
    (90:\RingRo) --
    (35:\RingRo) --
    (30:{\RingRo+\RingTip}) --
    (25:\RingRo) --
    (30:\RingRi) -- cycle;

  % --- Stage 2: Design (right, -30 to 30 deg) ---
  \fill[c2]
    (30:\RingRi) --
    (30:\RingRo) --
    (-25:\RingRo) --
    (-30:{\RingRo+\RingTip}) --
    (-35:\RingRo) --
    (-30:\RingRi) -- cycle;

  % --- Stage 3: Development (bottom-right, -90 to -30 deg) ---
  \fill[c3]
    (-30:\RingRi) --
    (-30:\RingRo) --
    (-85:\RingRo) --
    (-90:{\RingRo+\RingTip}) --
    (-95:\RingRo) --
    (-90:\RingRi) -- cycle;

  % --- Stage 4: Testing (bottom-left, -150 to -90 deg) ---
  \fill[c4]
    (-90:\RingRi) --
    (-90:\RingRo) --
    (-145:\RingRo) --
    (-150:{\RingRo+\RingTip}) --
    (-155:\RingRo) --
    (-150:\RingRi) -- cycle;

  % --- Stage 5: Deployment (left, 150 to 210 deg) ---
  \fill[c5]
    (-150:\RingRi) --
    (-150:\RingRo) --
    (155:\RingRo) --
    (150:{\RingRo+\RingTip}) --
    (145:\RingRo) --
    (150:\RingRi) -- cycle;

  % --- Stage 6: Review (top-left, 90 to 150 deg) ---
  \fill[c6]
    (150:\RingRi) --
    (150:\RingRo) --
    (95:\RingRo) --
    (90:{\RingRo+\RingTip}) --
    (85:\RingRo) --
    (90:\RingRi) -- cycle;

  %------------------------------------------------------------------------
  % Inner decorative rings (subtle outlines for polish).
  %------------------------------------------------------------------------
  \draw[kauprimary!30, line width=0.6pt] (0,0) circle ({\RingRi-0.15});
  \draw[kauprimary!20, line width=0.4pt] (0,0) circle ({\RingRo+\RingTip+0.15});

  %------------------------------------------------------------------------
  % Center title.
  %------------------------------------------------------------------------
  \node[center title] at (0,0.35) {\RL{منهجية}};
  \node[center title] at (0,-0.45) {\RL{أجايل}};
  \node[center subtitle] at (0,-1.15) {\RL{\EN{Agile Methodology}}};

  %------------------------------------------------------------------------
  % Stage labels (Arabic name + English gloss) outside each segment.
  % Position at the mid-angle of each 60-degree segment.
  %------------------------------------------------------------------------
  \node[stage label] at (60:\RingRL)  {\RL{التخطيط}\\[2pt]{\small\color{kausecondary}\EN{Planning}}};
  \node[stage label] at (0:\RingRL)   {\RL{التصميم}\\[2pt]{\small\color{kausecondary}\EN{Design}}};
  \node[stage label] at (-60:\RingRL) {\RL{التطوير}\\[2pt]{\small\color{kausecondary}\EN{Development}}};
  \node[stage label] at (-120:\RingRL){\RL{الاختبار}\\[2pt]{\small\color{kausecondary}\EN{Testing}}};
  \node[stage label] at (180:\RingRL) {\RL{الإطلاق}\\[2pt]{\small\color{kausecondary}\EN{Deployment}}};
  \node[stage label] at (120:\RingRL) {\RL{المراجعة}\\[2pt]{\small\color{kausecondary}\EN{Review}}};

  %------------------------------------------------------------------------
  % Numbered badges (1–6) on each segment, centered on the segment.
  %------------------------------------------------------------------------
  \node[badge] at (60:3.0)  {1};
  \node[badge] at (0:3.0)   {2};
  \node[badge] at (-60:3.0) {3};
  \node[badge] at (-120:3.0){4};
  \node[badge] at (180:3.0) {5};
  \node[badge] at (120:3.0) {6};

\end{tikzpicture}
}
%     \caption{...}
%     \label{...}
%   \end{figure}
%==============================================================================

% Radii (in cm):
%   Ri = inner radius (where chevron arrows start)
%   Ro = outer radius (where chevron arrows end)
%   Tip = chevron tip extension beyond Ro
%   RL  = label radius (outer ring text)
\def\RingRi{2.4}
\def\RingRo{3.6}
\def\RingTip{0.35}
\def\RingRL{4.95}

\begin{tikzpicture}[
  % Six stage color styles — drawn from the Bayan palette family.
  c1/.style={fill=kauprimary!85,   draw=kauprimary,   text=white},
  c2/.style={fill=kausecondary!85, draw=kausecondary, text=white},
  c3/.style={fill=kaudark!85,      draw=kaudark,      text=white},
  c4/.style={fill=kauaccent!85,    draw=kauaccent,    text=white},
  c5/.style={fill=kauprimary!55,   draw=kauprimary,   text=white},
  c6/.style={fill=kausecondary!50, draw=kausecondary, text=white},
  badge/.style={
    circle, fill=white, draw=kauprimary, line width=1pt,
    text=kauprimary, font=\large\bfseries, inner sep=0pt,
    minimum size=1.1cm,
  },
  stage label/.style={
    font=\large\bfseries, text=kaudark, align=center,
  },
  center title/.style={
    font=\Huge\bfseries\color{kauprimary}, align=center,
  },
  center subtitle/.style={
    font=\normalsize\color{kausecondary}, align=center,
  },
  >=stealth,
]

  %------------------------------------------------------------------------
  % Six chevron-arrow segments arranged in a clockwise ring.
  % Each segment spans 60 degrees. Math convention: 0°=right, 90°=top.
  % Stage 1: 30°→90° (top-right)        — Planning
  % Stage 2: -30°→30° (right)            — Design
  % Stage 3: -90°→-30° (bottom-right)    — Development
  % Stage 4: -150°→-90° (bottom-left)    — Testing
  % Stage 5: 150°→210° (left)            — Deployment  [210° = -150°]
  % Stage 6: 90°→150° (top-left)         — Review
  %------------------------------------------------------------------------

  % --- Stage 1: Planning (top-right, 30 to 90 deg) ---
  \fill[c1]
    (90:\RingRi) --
    (90:\RingRo) --
    (35:\RingRo) --
    (30:{\RingRo+\RingTip}) --
    (25:\RingRo) --
    (30:\RingRi) -- cycle;

  % --- Stage 2: Design (right, -30 to 30 deg) ---
  \fill[c2]
    (30:\RingRi) --
    (30:\RingRo) --
    (-25:\RingRo) --
    (-30:{\RingRo+\RingTip}) --
    (-35:\RingRo) --
    (-30:\RingRi) -- cycle;

  % --- Stage 3: Development (bottom-right, -90 to -30 deg) ---
  \fill[c3]
    (-30:\RingRi) --
    (-30:\RingRo) --
    (-85:\RingRo) --
    (-90:{\RingRo+\RingTip}) --
    (-95:\RingRo) --
    (-90:\RingRi) -- cycle;

  % --- Stage 4: Testing (bottom-left, -150 to -90 deg) ---
  \fill[c4]
    (-90:\RingRi) --
    (-90:\RingRo) --
    (-145:\RingRo) --
    (-150:{\RingRo+\RingTip}) --
    (-155:\RingRo) --
    (-150:\RingRi) -- cycle;

  % --- Stage 5: Deployment (left, 150 to 210 deg) ---
  \fill[c5]
    (-150:\RingRi) --
    (-150:\RingRo) --
    (155:\RingRo) --
    (150:{\RingRo+\RingTip}) --
    (145:\RingRo) --
    (150:\RingRi) -- cycle;

  % --- Stage 6: Review (top-left, 90 to 150 deg) ---
  \fill[c6]
    (150:\RingRi) --
    (150:\RingRo) --
    (95:\RingRo) --
    (90:{\RingRo+\RingTip}) --
    (85:\RingRo) --
    (90:\RingRi) -- cycle;

  %------------------------------------------------------------------------
  % Inner decorative rings (subtle outlines for polish).
  %------------------------------------------------------------------------
  \draw[kauprimary!30, line width=0.6pt] (0,0) circle ({\RingRi-0.15});
  \draw[kauprimary!20, line width=0.4pt] (0,0) circle ({\RingRo+\RingTip+0.15});

  %------------------------------------------------------------------------
  % Center title.
  %------------------------------------------------------------------------
  \node[center title] at (0,0.35) {\RL{منهجية}};
  \node[center title] at (0,-0.45) {\RL{أجايل}};
  \node[center subtitle] at (0,-1.15) {\RL{\EN{Agile Methodology}}};

  %------------------------------------------------------------------------
  % Stage labels (Arabic name + English gloss) outside each segment.
  % Position at the mid-angle of each 60-degree segment.
  %------------------------------------------------------------------------
  \node[stage label] at (60:\RingRL)  {\RL{التخطيط}\\[2pt]{\small\color{kausecondary}\EN{Planning}}};
  \node[stage label] at (0:\RingRL)   {\RL{التصميم}\\[2pt]{\small\color{kausecondary}\EN{Design}}};
  \node[stage label] at (-60:\RingRL) {\RL{التطوير}\\[2pt]{\small\color{kausecondary}\EN{Development}}};
  \node[stage label] at (-120:\RingRL){\RL{الاختبار}\\[2pt]{\small\color{kausecondary}\EN{Testing}}};
  \node[stage label] at (180:\RingRL) {\RL{الإطلاق}\\[2pt]{\small\color{kausecondary}\EN{Deployment}}};
  \node[stage label] at (120:\RingRL) {\RL{المراجعة}\\[2pt]{\small\color{kausecondary}\EN{Review}}};

  %------------------------------------------------------------------------
  % Numbered badges (1–6) on each segment, centered on the segment.
  %------------------------------------------------------------------------
  \node[badge] at (60:3.0)  {1};
  \node[badge] at (0:3.0)   {2};
  \node[badge] at (-60:3.0) {3};
  \node[badge] at (-120:3.0){4};
  \node[badge] at (180:3.0) {5};
  \node[badge] at (120:3.0) {6};

\end{tikzpicture}
}
  \caption{نموذج أجايل (\EN{Agile Model}) المتبع في مشروع بيان.}
  \label{fig:agile}
\end{figure}

تم اعتماد منهجية أجايل (\EN{Agile}) في تطوير مشروع بيان، نظراً لمرونتها
وقابليتها للتكيّف مع المتطلبات المتغيرة، خاصةً في مشروع يعتمد على تقنيات
ذكاء اصطناعي حديثة ومتطورة باستمرار. وتقوم المنهجية على تقسيم العمل إلى دورات
تطوير قصيرة (\EN{Iterations}) تمر كل دورة منها بست مراحل متتابعة: التخطيط
(\EN{Planning})، والتصميم (\EN{Design})، والتطوير (\EN{Development})،
والاختبار (\EN{Testing})، والإطلاق (\EN{Deployment})، والمراجعة
(\EN{Review})، مع التركيز على التسليم التدريجي والتغذية الراجعة المستمرة من
أصحاب المصلحة والمشرفين. ويوضح الشكل~\ref{fig:agile} نموذج أجايل المتبع في
المشروع، بينما تُفصَّل مراحل الدورة ومخرجاتها والدورات التنفيذية الأربع
للمشروع في قسم منهجية التطوير بالفصل الثالث
(القسم~\ref{sec:methodology}).

%------------------------------------------------------------------------------
\section{الجدول الزمني للمشروع (\EN{Project Schedule})}
\label{sec:schedule}

تُشتق الخطة الزمنية لمشروع «بيان» مباشرةً من هيكل تجزئة العمل
(\EN{Work Breakdown Structure}) المعتمد في هذا الفصل، وعلى ضوء المنهجية
الرشيقة (\EN{Agile}) الموصوفة في القسم~\ref{sec:methodology}؛ إذ حُوِّلت
الأنشطة المخطط لها إلى حزم عمل ومهام تنفيذية محددة المخرجات، ثم نُظّمت
زمنياً ضمن ثلاث مراحل رئيسية تتوافق مجتمعةً مع فصلي المشروعين الدراسيين
المقررين: المرحلة الأولى تختص بجهود التأسيس والتحليل والتوثيق (كتابة
الفصول الثلاثة الأولى) في الفصل الدراسي الأول، بينما تتوزع خطة التنفيذ
الكاملة (التصميم والتطوير والاختبار والتسليمات النهائية) على المرحلتين
الثانية والثالثة في الفصل الدراسي الثاني. ويمتد المشروع إجمالاً من \EN{01/07/2026}
حتى \EN{23/01/2027} بواقع 207 أيام تقويمية (نحو 30 أسبوعاً) تُحتسب
شاملةً نهايات الأسبوع، موزعةً على فصلين دراسيين بينهما فترة انتقالية تمتد
19 يوماً (من \EN{28/09/2026} إلى \EN{16/10/2026}) تُخصص للعطلة الفاصلة
بين الفصلين وتهيئة بيئة العمل للمرحلة الثانية.

واتُّخذ تقويم لجنة مشاريع التخرج للعام الجامعي \EN{2026/2027} مرجعاً
رسمياً في ضبط جميع نقاط التسليم المقررة؛ ابتداءً من إقرار المشاريع
وتعيين المشرفين، مروراً باستمارات المتابعة الدورية لتقارير الفصول، وصولاً
إلى موعدَي المناقشتين النهائيتين. وقد بُني الجدول باتجاه
عكسي انطلاقاً من موعدَي المناقشتين الثابتين (\EN{20/09/2026} و
\EN{16/01/2027})، بحيث تُشتق مواعيد بدء الأنشطة من مواعيد إنجازها
المستهدفة، مع الالتزام بمقتضى التقويم القاضي بتسليم نسختين من الوثيقة
النهائية قبل أسبوع على الأقل من المناقشة، إضافةً إلى البوستر العلمي
والفيديو التسويقي المطلوبَين لمشروع التخرج الثاني.

وتربط الخطة بين مصطلحات إدارة المشاريع ومفاهيم المنهجية الرشيقة على
النحو الآتي: تعمل حزم التوثيق في المرحلة الأولى وفق تدفق موجّه بالبوابات
(\EN{Gate-Driven}) تفرضه استمارات المتابعة الرسمية، بينما تقابل كل دورة
تطوير في المرحلة الثانية إصداراً وظيفياً متزايداً (\EN{Increment}) تُعرض
مخرجاته على المشرف في مراجعة نهاية الدورة قبل اعتماده ضمن النظام. ويعرض
هذا القسم الخطة الزمنية عبر مستويين متكاملين يُغني كلٌّ منهما عن الآخر:
مستوى الملخص التنفيذي الذي يجسده الجدول المدمج للمراحل الكبرى
ومعالمها الرسمية (الجدول~\ref{tab:unified_schedule})،
ومستوى التنفيذ التفصيلي الذي يجسده مخططا جانت الشاملان بكامل المهام
الثمانين والثلاثة موزعةً على صفوفهما على نمط لوحة العرض المعتمد في برامج
إدارة المشاريع الاحترافية: لوحة بيانات جانبية لكل صف تثبت رمز المهمة
ومددتها وتاريخَيها ومسؤولها بجوار شريطها الزمني الملوَّن وفق المسؤول عنها
(القسم الفرعي~\ref{subsec:gantt})؛ فلا يتكرر عرض أي
معلومة في المستويين، ويظل الجدول الملخص مرجع الحوكمة الزمنية العليا
بينما يظل المخططان مرجع المتابعة الأسبوعية التفصيلية.

%------------------------------------------------------------------------------
% الجدول الملخص المدمج للمراحل والمعالم: يحل محل القسم الفرعي السابق
% (خط الأساس الزمني المُلخَّص وسجل المعالم) وجداوله الثلاثة
تمتد الخطة الزمنية لمشروع «بيان» على ثلاث مراحل رئيسية تتوافق مجتمعةً مع
فصلي مشروع التخرج للعام الجامعي \EN{2026/2027}. وبناءً على منهجية «أجايل» المعتمدة،
قُسمت المهام التقنية إلى دورات تطويرية (\EN{Sprints}) تركز على التسليم
التدريجي، بالتوازي مع مهام التوثيق الأكاديمي. ويلخص
الجدول~\ref{tab:unified_schedule} المراحل الكبرى للمشروع ومعالمه الرسمية،
بينما تُعرض التفاصيل التنفيذية وتوزيع مهام الفريق في مخططات جانت
(الشكل~\ref{fig:ganttP1} والشكل~\ref{fig:ganttP2}).

\begin{table}[H]
\centering
\caption{الملخص الزمني لمراحل تنفيذ مشروع بيان والمعالم الرئيسية}
\label{tab:unified_schedule}
\small
\setlength{\tabcolsep}{4pt}
\renewcommand{\arraystretch}{1.5}
\begin{tabular}{|p{3.5cm}|p{2.5cm}|p{5cm}|p{3.5cm}|}
\hline
\textbf{المرحلة الرئيسية} & \textbf{الفترة الزمنية} & \textbf{أبرز الأنشطة المنجزة (وفق أجايل)} & \textbf{أهم المعالم (\EN{Milestones})} \\
\hline
\textbf{المرحلة الأولى: التحليل والتأسيس} (مشروع تخرج 1) & يوليو -- سبتمبر 2026 & جمع المتطلبات، تحليل الجدوى، وتصميم معمارية النظام (الفصول 1، 2، 3). & تسليم وثيقة مشروع تخرج 1، واجتياز المناقشة الأولى (\EN{M1}--\EN{M6}). \\
\hline
\textbf{المرحلة الثانية: التطوير التقني} (مشروع تخرج 2) & أكتوبر -- ديسمبر 2026 & تنفيذ 4 دورات تطوير (\EN{Sprints}) لبناء الوكلاء (\EN{SQL}، \EN{RAG})، دمج النظام، وتطوير الواجهات. & اعتماد تقارير الإنجاز الدورية، وتقديم نسخة برمجية قابلة للتشغيل (\EN{M8}--\EN{M10}). \\
\hline
\textbf{المرحلة الثالثة: الاختبار والتسليم} & يناير 2027 & الاختبار الشامل (\EN{UAT})، التوثيق النهائي (الفصول 4--7)، وتجهيز العرض والبوستر. & تسليم الوثيقة النهائية، واجتياز المناقشة الختامية (\EN{M11}--\EN{M13}). \\
\hline
\end{tabular}
\end{table}

%------------------------------------------------------------------------------
\subsection{مخططا جانت التنفيذيان (\EN{Detailed Gantt Charts})}
\label{subsec:gantt}

تنزَّل الخطة الزمنية إلى مستوى التنفيذ اليومي في مخططَي جانت الشاملين
بالشكلين~\ref{fig:ganttP1} و\ref{fig:ganttP2}؛ إذ يعرض الشكل~\ref{fig:ganttP1}
المرحلة الأولى كاملة (يوليو -- سبتمبر 2026) بجميع مهامها الأربع والأربعين
موزعةً على صفوف حزمها الخمس، بينما يعرض الشكل~\ref{fig:ganttP2} المرحلتين
الثانية والثالثة معاً (أكتوبر 2026 -- يناير 2027) بمهامها التسع والثلاثين
موزعةً على صفوف حزمها الست، بعد أن كانت الفترة الانتقالية بين الفصلين
موثقةً في
مقدمة هذا القسم. وقد صُمِّم المخططان على نمط لوحة العرض المعتمد في
برامج إدارة المشاريع الاحترافية (مثل \EN{MS Project}): لوحة واحدة
موحَّدة بإطار مشترك وشبكة صفوف ممتدة من أقصى اليمين إلى أقصى اليسار،
تجمع لوحة بيانات المهام على يمينها والخط الزمني الرسومي على يسارها في
المسح البصري نفسه، بحيث يحمل كل صف مهمةً واحدة تظهر بياناتها الكاملة
في أقصى اليمين ويبدأ شريطها الزمني ملاصقاً لها متجهاً نحو اليسار.

تضم لوحة بيانات المهام ستة أعمدة مثبتة لكل صف: رمز المهمة في هيكل تجزئة
العمل (\EN{WBS})، واسمها التنفيذي، ومدتها بالأيام التقويمية، وتاريخَي
بدايتها ونهايتها بصيغة اليوم/الشهر، والمسؤول عنها برمز تخصصه
(\EN{AI} للذكاء الاصطناعي، و\EN{BE} للواجهة الخلفية وقواعد البيانات،
و\EN{FE} للواجهة الأمامية وتجربة المستخدم، و\EN{TM} للفريق كاملاً)، مع
صفوف الحزم المظللة التي تحمل رؤوساً بارزة تُجمِّع مهامها وتثبت حدود
الحزمة في خط الأساس المعتمد في الجدول~\ref{tab:unified_schedule}. أما الجدول الزمني
الرسومي فيعرض أشرطة المهام الملوَّنة وفق المسؤول عنها، وأشرطة الحزم
الملخَّصة الرفيعة الداكنة التي تجسد حدود كل حزمة من أول مهمة فيها إلى
آخرها، ومعينات المعالم الرسمية الحمراء المُضمَّنة في صفوف المهام التي
تُنجز عندها ملحقةً بتاريخها الرسمي، وأسهم التبعية التي تربط عناصر المسار
الحرج. وتُظلَّل خلايا الجمعة والسبت بظلٍّ رمادي يميز عطلة نهاية الأسبوع، بينما
تتقدم التقويمة الميلادية من يمين المحور الملاصق للوحة البيانات إلى
يساره بمقدار خانة واحدة لكل يوم، تحت ترويسة شهرية وأسابيع مرقمة تبدأ
بيوم الأحد، مع خط عمودي رفيع عند بداية كل أسبوع ييسّر تتبع التقدم
ومقارنة المخطط بالإنجاز الفعلي في اجتماعات المتابعة.

\begin{figure}[H]
  \centering
\definecolor{ownerAI}{HTML}{0E7490}%
\definecolor{ownerBE}{HTML}{2E75B6}%
\definecolor{ownerFE}{HTML}{6D28D9}%
\definecolor{ownerTM}{HTML}{64748B}%
\newsavebox\gCBA
\begin{lrbox}{\gCBA}%
\begin{LTR}%
\begin{tikzpicture}[x=1pt, y=1pt]
\fill[kauprimary] (0,523.2) rectangle (517.8,535.9);
\fill[kauprimary!14] (0,511.5) rectangle (293.8,523.2);
\fill[kauprimary!7] (0,501.3) rectangle (293.8,511.5);
\fill[kauprimary!22] (293.8,501.3) rectangle (517.8,523.2);
\fill[kauprimary!10] (0,491.0) rectangle (517.8,501.3);
\fill[kauprimary!10] (0,429.7) rectangle (517.8,439.9);
\fill[kauprimary!10] (0,337.6) rectangle (517.8,347.8);
\fill[kauprimary!10] (0,214.8) rectangle (517.8,225.1);
\fill[kauprimary!10] (0,61.4) rectangle (517.8,71.6);
\fill[gray!38] (280.6,511.5) rectangle (287.2,0.0);
\fill[gray!38] (257.5,511.5) rectangle (264.1,0.0);
\fill[gray!38] (234.4,511.5) rectangle (241.0,0.0);
\fill[gray!38] (211.3,511.5) rectangle (217.9,0.0);
\fill[gray!38] (188.2,511.5) rectangle (194.8,0.0);
\fill[gray!38] (165.1,511.5) rectangle (171.7,0.0);
\fill[gray!38] (142.0,511.5) rectangle (148.6,0.0);
\fill[gray!38] (118.9,511.5) rectangle (125.5,0.0);
\fill[gray!38] (95.7,511.5) rectangle (102.3,0.0);
\fill[gray!38] (72.6,511.5) rectangle (79.2,0.0);
\fill[gray!38] (49.5,511.5) rectangle (56.1,0.0);
\fill[gray!38] (26.4,511.5) rectangle (33.0,0.0);
\fill[gray!38] (3.3,511.5) rectangle (9.9,0.0);
\draw[gray!55, line width=0.45pt] (191.5,0) -- (191.5,523.2);
\draw[gray!55, line width=0.45pt] (89.1,0) -- (89.1,523.2);
\draw[gray!45, line width=0.3pt] (280.6,0) -- (280.6,511.5);
\draw[gray!45, line width=0.3pt] (257.5,0) -- (257.5,511.5);
\draw[gray!45, line width=0.3pt] (234.4,0) -- (234.4,511.5);
\draw[gray!45, line width=0.3pt] (211.3,0) -- (211.3,511.5);
\draw[gray!45, line width=0.3pt] (188.2,0) -- (188.2,511.5);
\draw[gray!45, line width=0.3pt] (165.1,0) -- (165.1,511.5);
\draw[gray!45, line width=0.3pt] (142.0,0) -- (142.0,511.5);
\draw[gray!45, line width=0.3pt] (118.9,0) -- (118.9,511.5);
\draw[gray!45, line width=0.3pt] (95.7,0) -- (95.7,511.5);
\draw[gray!45, line width=0.3pt] (72.6,0) -- (72.6,511.5);
\draw[gray!45, line width=0.3pt] (49.5,0) -- (49.5,511.5);
\draw[gray!45, line width=0.3pt] (26.4,0) -- (26.4,511.5);
\draw[gray!45, line width=0.3pt] (3.3,0) -- (3.3,511.5);
\fill[kaudark!90] (257.5,494.7) rectangle (293.8,497.0);
\filldraw[fill=ownerTM!50, draw=ownerTM, line width=0.4pt] (283.9,482.9) rectangle (293.8,489.0);
\filldraw[fill=ownerTM!50, draw=ownerTM, line width=0.4pt] (277.3,472.6) rectangle (283.9,478.8);
\filldraw[fill=ownerBE!55, draw=ownerBE, line width=0.4pt] (267.4,462.4) rectangle (280.6,468.5);
\filldraw[fill=ownerTM!50, draw=ownerTM, line width=0.4pt] (260.8,452.2) rectangle (270.7,458.3);
\filldraw[fill=ownerTM!50, draw=ownerTM, line width=0.4pt] (257.5,441.9) rectangle (264.1,448.1);
\fill[kaudark!90] (194.8,433.3) rectangle (293.8,435.6);
\filldraw[fill=ownerAI!55, draw=ownerAI, line width=0.4pt] (260.8,421.5) rectangle (293.8,427.6);
\filldraw[fill=ownerBE!55, draw=ownerBE, line width=0.4pt] (244.3,411.2) rectangle (270.7,417.4);
\filldraw[fill=ownerTM!50, draw=ownerTM, line width=0.4pt] (234.4,401.0) rectangle (247.6,407.2);
\filldraw[fill=ownerTM!50, draw=ownerTM, line width=0.4pt] (224.5,390.8) rectangle (237.7,396.9);
\filldraw[fill=ownerFE!45, draw=ownerFE, line width=0.4pt] (217.9,380.6) rectangle (227.8,386.7);
\filldraw[fill=ownerAI!55, draw=ownerAI, line width=0.4pt] (211.3,370.3) rectangle (221.2,376.5);
\filldraw[fill=ownerTM!50, draw=ownerTM, line width=0.4pt] (201.4,360.1) rectangle (214.6,366.2);
\filldraw[fill=ownerTM!50, draw=ownerTM, line width=0.4pt] (194.8,349.9) rectangle (204.7,356.0);
\fill[kaudark!90] (125.5,341.3) rectangle (194.8,343.5);
\filldraw[fill=ownerAI!55, draw=ownerAI, line width=0.4pt] (184.9,329.4) rectangle (194.8,335.5);
\filldraw[fill=ownerAI!55, draw=ownerAI, line width=0.4pt] (175.0,319.2) rectangle (184.9,325.3);
\filldraw[fill=ownerBE!55, draw=ownerBE, line width=0.4pt] (165.1,308.9) rectangle (178.3,315.1);
\filldraw[fill=ownerAI!55, draw=ownerAI, line width=0.4pt] (158.5,298.7) rectangle (168.4,304.9);
\filldraw[fill=ownerAI!55, draw=ownerAI, line width=0.4pt] (151.9,288.5) rectangle (161.8,294.6);
\filldraw[fill=ownerAI!55, draw=ownerAI, line width=0.4pt] (145.3,278.3) rectangle (155.2,284.4);
\filldraw[fill=ownerFE!45, draw=ownerFE, line width=0.4pt] (138.7,268.0) rectangle (151.9,274.2);
\filldraw[fill=ownerTM!50, draw=ownerTM, line width=0.4pt] (135.4,257.8) rectangle (142.0,263.9);
\filldraw[fill=ownerAI!55, draw=ownerAI, line width=0.4pt] (132.1,247.6) rectangle (138.7,253.7);
\filldraw[fill=ownerAI!55, draw=ownerAI, line width=0.4pt] (128.8,237.3) rectangle (135.4,243.5);
\filldraw[fill=ownerTM!50, draw=ownerTM, line width=0.4pt] (125.5,227.1) rectangle (132.1,233.2);
\fill[kaudark!90] (46.2,218.5) rectangle (125.5,220.8);
\filldraw[fill=ownerTM!50, draw=ownerTM, line width=0.4pt] (115.6,206.6) rectangle (125.5,212.8);
\filldraw[fill=ownerTM!50, draw=ownerTM, line width=0.4pt] (112.3,196.4) rectangle (118.9,202.6);
\filldraw[fill=ownerBE!55, draw=ownerBE, line width=0.4pt] (102.3,186.2) rectangle (115.6,192.3);
\filldraw[fill=ownerBE!55, draw=ownerBE, line width=0.4pt] (92.4,176.0) rectangle (105.6,182.1);
\filldraw[fill=ownerTM!50, draw=ownerTM, line width=0.4pt] (85.8,165.7) rectangle (95.7,171.9);
\filldraw[fill=ownerBE!55, draw=ownerBE, line width=0.4pt] (79.2,155.5) rectangle (89.1,161.6);
\filldraw[fill=ownerBE!55, draw=ownerBE, line width=0.4pt] (72.6,145.3) rectangle (82.5,151.4);
\filldraw[fill=ownerBE!55, draw=ownerBE, line width=0.4pt] (69.3,135.0) rectangle (75.9,141.2);
\filldraw[fill=ownerFE!45, draw=ownerFE, line width=0.4pt] (62.7,124.8) rectangle (72.6,130.9);
\filldraw[fill=ownerFE!45, draw=ownerFE, line width=0.4pt] (56.1,114.6) rectangle (66.0,120.7);
\filldraw[fill=ownerBE!55, draw=ownerBE, line width=0.4pt] (52.8,104.3) rectangle (59.4,110.5);
\filldraw[fill=ownerTM!50, draw=ownerTM, line width=0.4pt] (49.5,94.1) rectangle (56.1,100.3);
\filldraw[fill=ownerTM!50, draw=ownerTM, line width=0.4pt] (46.2,83.9) rectangle (52.8,90.0);
\filldraw[fill=ownerTM!50, draw=ownerTM, line width=0.4pt] (46.2,73.7) rectangle (52.8,79.8);
\fill[kaudark!90] (0.0,65.1) rectangle (89.1,67.3);
\filldraw[fill=ownerTM!50, draw=ownerTM, line width=0.4pt] (69.3,53.2) rectangle (89.1,59.3);
\filldraw[fill=ownerTM!50, draw=ownerTM, line width=0.4pt] (56.1,43.0) rectangle (75.9,49.1);
\filldraw[fill=ownerTM!50, draw=ownerTM, line width=0.4pt] (46.2,32.7) rectangle (56.1,38.9);
\filldraw[fill=kauaccent!35, draw=kauaccent!70!black, line width=0.4pt] (36.3,22.5) rectangle (46.2,28.6);
\filldraw[fill=kauaccent!35, draw=kauaccent!70!black, line width=0.4pt] (26.4,12.3) rectangle (36.3,18.4);
\filldraw[fill=kauaccent!75, draw=kauaccent!80!black, line width=0.4pt] (0.0,2.0) rectangle (26.4,8.2);
\node[diamond, fill=kauaccent, draw=kauaccent!60, inner sep=2.6pt, line width=0.4pt] at (275.7,475.7) {};
\node[anchor=east, font=\tiny\bfseries, inner sep=0pt] at (271.5,475.7) {\EN{M1 06/07}};
\node[diamond, fill=kauaccent, draw=kauaccent!60, inner sep=2.6pt, line width=0.4pt] at (259.2,465.5) {};
\node[anchor=east, font=\tiny\bfseries, inner sep=0pt] at (255.0,465.5) {\EN{M2 11/07}};
\node[diamond, fill=kauaccent, draw=kauaccent!60, inner sep=2.6pt, line width=0.4pt] at (196.4,352.9) {};
\node[anchor=east, font=\tiny\bfseries, inner sep=0pt] at (192.2,352.9) {\EN{M3 30/07}};
\node[diamond, fill=kauaccent, draw=kauaccent!60, inner sep=2.6pt, line width=0.4pt] at (127.1,230.2) {};
\node[anchor=east, font=\tiny\bfseries, inner sep=0pt] at (122.9,230.2) {\EN{M4 20/08}};
\node[diamond, fill=kauaccent, draw=kauaccent!60, inner sep=2.6pt, line width=0.4pt] at (47.9,76.7) {};
\node[anchor=south, font=\tiny\bfseries, inner sep=0pt] at (47.9,79.2) {\EN{M5 13/09}};
\node[diamond, fill=kauaccent, draw=kauaccent!60, inner sep=2.6pt, line width=0.4pt] at (24.8,15.3) {};
\node[anchor=south, font=\tiny\bfseries, inner sep=0pt] at (24.8,17.8) {\EN{M6 20/09}};
\draw[-stealth, kauprimary, line width=0.75pt] (194.8,434.8) -- (191.8,434.8) -- (191.8,342.7) -- (125.5,342.7);
\draw[-stealth, kauprimary, line width=0.75pt] (125.5,342.7) -- (122.5,342.7) -- (122.5,219.9) -- (46.2,219.9);
\draw[-stealth, kauprimary, line width=0.75pt] (46.2,219.9) -- (43.2,219.9) -- (43.2,66.5) -- (0.0,66.5);
\draw[kauprimary!50, line width=0.5pt] (0,501.3) rectangle (517.8,501.3);
\draw[kauprimary!45, line width=0.45pt] (0,491.0) -- (517.8,491.0);
\draw[gray!32, line width=0.3pt] (0,480.8) -- (517.8,480.8);
\draw[gray!32, line width=0.3pt] (0,470.6) -- (517.8,470.6);
\draw[gray!32, line width=0.3pt] (0,460.4) -- (517.8,460.4);
\draw[gray!32, line width=0.3pt] (0,450.1) -- (517.8,450.1);
\draw[gray!32, line width=0.3pt] (0,439.9) -- (517.8,439.9);
\draw[kauprimary!45, line width=0.45pt] (0,429.7) -- (517.8,429.7);
\draw[gray!32, line width=0.3pt] (0,419.4) -- (517.8,419.4);
\draw[gray!32, line width=0.3pt] (0,409.2) -- (517.8,409.2);
\draw[gray!32, line width=0.3pt] (0,399.0) -- (517.8,399.0);
\draw[gray!32, line width=0.3pt] (0,388.7) -- (517.8,388.7);
\draw[gray!32, line width=0.3pt] (0,378.5) -- (517.8,378.5);
\draw[gray!32, line width=0.3pt] (0,368.3) -- (517.8,368.3);
\draw[gray!32, line width=0.3pt] (0,358.1) -- (517.8,358.1);
\draw[gray!32, line width=0.3pt] (0,347.8) -- (517.8,347.8);
\draw[kauprimary!45, line width=0.45pt] (0,337.6) -- (517.8,337.6);
\draw[gray!32, line width=0.3pt] (0,327.4) -- (517.8,327.4);
\draw[gray!32, line width=0.3pt] (0,317.1) -- (517.8,317.1);
\draw[gray!32, line width=0.3pt] (0,306.9) -- (517.8,306.9);
\draw[gray!32, line width=0.3pt] (0,296.7) -- (517.8,296.7);
\draw[gray!32, line width=0.3pt] (0,286.4) -- (517.8,286.4);
\draw[gray!32, line width=0.3pt] (0,276.2) -- (517.8,276.2);
\draw[gray!32, line width=0.3pt] (0,266.0) -- (517.8,266.0);
\draw[gray!32, line width=0.3pt] (0,255.8) -- (517.8,255.8);
\draw[gray!32, line width=0.3pt] (0,245.5) -- (517.8,245.5);
\draw[gray!32, line width=0.3pt] (0,235.3) -- (517.8,235.3);
\draw[gray!32, line width=0.3pt] (0,225.1) -- (517.8,225.1);
\draw[kauprimary!45, line width=0.45pt] (0,214.8) -- (517.8,214.8);
\draw[gray!32, line width=0.3pt] (0,204.6) -- (517.8,204.6);
\draw[gray!32, line width=0.3pt] (0,194.4) -- (517.8,194.4);
\draw[gray!32, line width=0.3pt] (0,184.1) -- (517.8,184.1);
\draw[gray!32, line width=0.3pt] (0,173.9) -- (517.8,173.9);
\draw[gray!32, line width=0.3pt] (0,163.7) -- (517.8,163.7);
\draw[gray!32, line width=0.3pt] (0,153.5) -- (517.8,153.5);
\draw[gray!32, line width=0.3pt] (0,143.2) -- (517.8,143.2);
\draw[gray!32, line width=0.3pt] (0,133.0) -- (517.8,133.0);
\draw[gray!32, line width=0.3pt] (0,122.8) -- (517.8,122.8);
\draw[gray!32, line width=0.3pt] (0,112.5) -- (517.8,112.5);
\draw[gray!32, line width=0.3pt] (0,102.3) -- (517.8,102.3);
\draw[gray!32, line width=0.3pt] (0,92.1) -- (517.8,92.1);
\draw[gray!32, line width=0.3pt] (0,81.8) -- (517.8,81.8);
\draw[gray!32, line width=0.3pt] (0,71.6) -- (517.8,71.6);
\draw[kauprimary!45, line width=0.45pt] (0,61.4) -- (517.8,61.4);
\draw[gray!32, line width=0.3pt] (0,51.2) -- (517.8,51.2);
\draw[gray!32, line width=0.3pt] (0,40.9) -- (517.8,40.9);
\draw[gray!32, line width=0.3pt] (0,30.7) -- (517.8,30.7);
\draw[gray!32, line width=0.3pt] (0,20.5) -- (517.8,20.5);
\draw[gray!32, line width=0.3pt] (0,10.2) -- (517.8,10.2);
\draw[gray!32, line width=0.3pt] (0,0.0) -- (517.8,0.0);
\draw[kauprimary!70, line width=0.6pt] (0,523.2) -- (517.8,523.2);
\draw[kauprimary!35, line width=0.4pt] (0,511.5) -- (293.8,511.5);
\draw[gray!35, line width=0.3pt] (315.8,0) -- (315.8,501.3);
\draw[kauprimary!55, line width=0.4pt] (315.8,501.3) -- (315.8,523.2);
\draw[gray!35, line width=0.3pt] (440.8,0) -- (440.8,501.3);
\draw[kauprimary!55, line width=0.4pt] (440.8,501.3) -- (440.8,523.2);
\draw[gray!35, line width=0.3pt] (457.8,0) -- (457.8,501.3);
\draw[kauprimary!55, line width=0.4pt] (457.8,501.3) -- (457.8,523.2);
\draw[gray!35, line width=0.3pt] (477.3,0) -- (477.3,501.3);
\draw[kauprimary!55, line width=0.4pt] (477.3,501.3) -- (477.3,523.2);
\draw[gray!35, line width=0.3pt] (496.8,0) -- (496.8,501.3);
\draw[kauprimary!55, line width=0.4pt] (496.8,501.3) -- (496.8,523.2);
\draw[kauprimary!60, line width=0.65pt] (293.8,0) -- (293.8,523.2);
\draw[kauprimary!70, line width=0.8pt] (0,0) rectangle (517.8,535.9);
\node[font=\scriptsize\bfseries, text=white, inner sep=0pt] at (258.9,529.5) {\textarabic{المرحلة الأولى — مشروع التخرج 1: التأسيس والتحليل والتوثيق}};
\node[font=\tiny\bfseries, inner sep=0pt] at (242.7,517.3) {\textarabic{يوليو 2026}};
\node[font=\tiny\bfseries, inner sep=0pt] at (140.3,517.3) {\textarabic{أغسطس}};
\node[font=\tiny\bfseries, inner sep=0pt] at (44.6,517.3) {\textarabic{سبتمبر}};
\node[font=\tiny, inner sep=0pt] at (269.1,506.4) {\textarabic{١}};
\node[font=\tiny, inner sep=0pt] at (246.0,506.4) {\textarabic{٢}};
\node[font=\tiny, inner sep=0pt] at (222.9,506.4) {\textarabic{٣}};
\node[font=\tiny, inner sep=0pt] at (199.7,506.4) {\textarabic{٤}};
\node[font=\tiny, inner sep=0pt] at (176.6,506.4) {\textarabic{٥}};
\node[font=\tiny, inner sep=0pt] at (153.5,506.4) {\textarabic{٦}};
\node[font=\tiny, inner sep=0pt] at (130.4,506.4) {\textarabic{٧}};
\node[font=\tiny, inner sep=0pt] at (107.3,506.4) {\textarabic{٨}};
\node[font=\tiny, inner sep=0pt] at (84.2,506.4) {\textarabic{٩}};
\node[font=\tiny, inner sep=0pt] at (61.1,506.4) {\textarabic{١٠}};
\node[font=\tiny, inner sep=0pt] at (38.0,506.4) {\textarabic{١١}};
\node[font=\tiny, inner sep=0pt] at (14.9,506.4) {\textarabic{١٢}};
\node[font=\tiny, inner sep=0pt] at (1.7,506.4) {\textarabic{١٣}};
\node[font=\tiny\bfseries, inner sep=0pt] at (304.8,512.2) {\textarabic{الرمز}};
\node[font=\tiny\bfseries, inner sep=0pt] at (378.3,512.2) {\textarabic{المهمة}};
\node[font=\tiny\bfseries, inner sep=0pt] at (449.3,512.2) {\textarabic{الأيام}};
\node[font=\tiny\bfseries, inner sep=0pt] at (467.6,512.2) {\textarabic{البداية}};
\node[font=\tiny\bfseries, inner sep=0pt] at (487.1,512.2) {\textarabic{النهاية}};
\node[font=\tiny\bfseries, inner sep=0pt] at (507.3,512.2) {\textarabic{المسؤول}};
\node[font=\tiny, inner sep=0pt] at (304.8,496.2) {\textarabic{١.١}};
\node[font=\tiny\bfseries, inner sep=0pt] at (378.3,496.2) {\textbf{\textarabic{حزمة التهيئة والتخطيط}}};
\node[font=\tiny, inner sep=0pt] at (449.3,496.2) {\textarabic{١١}};
\node[font=\tiny, inner sep=0pt] at (467.6,496.2) {\EN{01/07}};
\node[font=\tiny, inner sep=0pt] at (487.1,496.2) {\EN{11/07}};
\node[font=\tiny, inner sep=0pt] at (507.3,496.2) {\textarabic{الفريق}};
\node[font=\tiny, inner sep=0pt] at (304.8,485.9) {\textarabic{١.١.١}};
\node[font=\scriptsize, inner sep=0pt] at (378.3,485.9) {\textarabic{تشكيل الفريق وتوزيع الأدوار}};
\node[font=\tiny, inner sep=0pt] at (449.3,485.9) {\textarabic{٣}};
\node[font=\tiny, inner sep=0pt] at (467.6,485.9) {\EN{01/07}};
\node[font=\tiny, inner sep=0pt] at (487.1,485.9) {\EN{03/07}};
\node[font=\tiny, inner sep=0pt] at (507.3,485.9) {\EN{TM}};
\node[font=\tiny, inner sep=0pt] at (304.8,475.7) {\textarabic{١.١.٢}};
\node[font=\scriptsize, inner sep=0pt] at (378.3,475.7) {\textarabic{الاجتماع التأسيسي وآلية المراجعات}};
\node[font=\tiny, inner sep=0pt] at (449.3,475.7) {\textarabic{٢}};
\node[font=\tiny, inner sep=0pt] at (467.6,475.7) {\EN{04/07}};
\node[font=\tiny, inner sep=0pt] at (487.1,475.7) {\EN{05/07}};
\node[font=\tiny, inner sep=0pt] at (507.3,475.7) {\EN{TM}};
\node[font=\tiny, inner sep=0pt] at (304.8,465.5) {\textarabic{١.١.٣}};
\node[font=\scriptsize, inner sep=0pt] at (378.3,465.5) {\textarabic{مساحة العمل الرقمية والمستودعات}};
\node[font=\tiny, inner sep=0pt] at (449.3,465.5) {\textarabic{٤}};
\node[font=\tiny, inner sep=0pt] at (467.6,465.5) {\EN{05/07}};
\node[font=\tiny, inner sep=0pt] at (487.1,465.5) {\EN{08/07}};
\node[font=\tiny, inner sep=0pt] at (507.3,465.5) {\EN{BE}};
\node[font=\tiny, inner sep=0pt] at (304.8,455.2) {\textarabic{١.١.٤}};
\node[font=\scriptsize, inner sep=0pt] at (378.3,455.2) {\textarabic{أدوات الإدارة والتوثيق والتصميم}};
\node[font=\tiny, inner sep=0pt] at (449.3,455.2) {\textarabic{٣}};
\node[font=\tiny, inner sep=0pt] at (467.6,455.2) {\EN{08/07}};
\node[font=\tiny, inner sep=0pt] at (487.1,455.2) {\EN{10/07}};
\node[font=\tiny, inner sep=0pt] at (507.3,455.2) {\EN{TM}};
\node[font=\tiny, inner sep=0pt] at (304.8,445.0) {\textarabic{١.١.٥}};
\node[font=\scriptsize, inner sep=0pt] at (378.3,445.0) {\textarabic{مسودة النطاق وحدود المشروع}};
\node[font=\tiny, inner sep=0pt] at (449.3,445.0) {\textarabic{٢}};
\node[font=\tiny, inner sep=0pt] at (467.6,445.0) {\EN{10/07}};
\node[font=\tiny, inner sep=0pt] at (487.1,445.0) {\EN{11/07}};
\node[font=\tiny, inner sep=0pt] at (507.3,445.0) {\EN{TM}};
\node[font=\tiny, inner sep=0pt] at (304.8,434.8) {\textarabic{١.٢}};
\node[font=\tiny\bfseries, inner sep=0pt] at (378.3,434.8) {\textbf{\textarabic{حزمة الفصل الأول — التأطير والمقدمة}}};
\node[font=\tiny, inner sep=0pt] at (449.3,434.8) {\textarabic{٣٠}};
\node[font=\tiny, inner sep=0pt] at (467.6,434.8) {\EN{01/07}};
\node[font=\tiny, inner sep=0pt] at (487.1,434.8) {\EN{30/07}};
\node[font=\tiny, inner sep=0pt] at (507.3,434.8) {\textarabic{الفريق}};
\node[font=\tiny, inner sep=0pt] at (304.8,424.5) {\textarabic{١.٢.١}};
\node[font=\scriptsize, inner sep=0pt] at (378.3,424.5) {\textarabic{جمع التقارير المرجعية وتحليلها}};
\node[font=\tiny, inner sep=0pt] at (449.3,424.5) {\textarabic{١٠}};
\node[font=\tiny, inner sep=0pt] at (467.6,424.5) {\EN{01/07}};
\node[font=\tiny, inner sep=0pt] at (487.1,424.5) {\EN{10/07}};
\node[font=\tiny, inner sep=0pt] at (507.3,424.5) {\EN{AI}};
\node[font=\tiny, inner sep=0pt] at (304.8,414.3) {\textarabic{١.٢.٢}};
\node[font=\scriptsize, inner sep=0pt] at (378.3,414.3) {\textarabic{دراسة السياق المؤسسي المحلي}};
\node[font=\tiny, inner sep=0pt] at (449.3,414.3) {\textarabic{٨}};
\node[font=\tiny, inner sep=0pt] at (467.6,414.3) {\EN{08/07}};
\node[font=\tiny, inner sep=0pt] at (487.1,414.3) {\EN{15/07}};
\node[font=\tiny, inner sep=0pt] at (507.3,414.3) {\EN{BE}};
\node[font=\tiny, inner sep=0pt] at (304.8,404.1) {\textarabic{١.٢.٣}};
\node[font=\scriptsize, inner sep=0pt] at (378.3,404.1) {\textarabic{صياغة بيان المشكلة ومظاهرها}};
\node[font=\tiny, inner sep=0pt] at (449.3,404.1) {\textarabic{٤}};
\node[font=\tiny, inner sep=0pt] at (467.6,404.1) {\EN{15/07}};
\node[font=\tiny, inner sep=0pt] at (487.1,404.1) {\EN{18/07}};
\node[font=\tiny, inner sep=0pt] at (507.3,404.1) {\EN{TM}};
\node[font=\tiny, inner sep=0pt] at (304.8,393.9) {\textarabic{١.٢.٤}};
\node[font=\scriptsize, inner sep=0pt] at (378.3,393.9) {\textarabic{تحديد الأهداف السبعة للمشروع}};
\node[font=\tiny, inner sep=0pt] at (449.3,393.9) {\textarabic{٤}};
\node[font=\tiny, inner sep=0pt] at (467.6,393.9) {\EN{18/07}};
\node[font=\tiny, inner sep=0pt] at (487.1,393.9) {\EN{21/07}};
\node[font=\tiny, inner sep=0pt] at (507.3,393.9) {\EN{TM}};
\node[font=\tiny, inner sep=0pt] at (304.8,383.6) {\textarabic{١.٢.٥}};
\node[font=\scriptsize, inner sep=0pt] at (378.3,383.6) {\textarabic{صياغة أهمية المشروع}};
\node[font=\tiny, inner sep=0pt] at (449.3,383.6) {\textarabic{٣}};
\node[font=\tiny, inner sep=0pt] at (467.6,383.6) {\EN{21/07}};
\node[font=\tiny, inner sep=0pt] at (487.1,383.6) {\EN{23/07}};
\node[font=\tiny, inner sep=0pt] at (507.3,383.6) {\EN{FE}};
\node[font=\tiny, inner sep=0pt] at (304.8,373.4) {\textarabic{١.٢.٦}};
\node[font=\scriptsize, inner sep=0pt] at (378.3,373.4) {\textarabic{توثيق مساهمات المشروع الأربع}};
\node[font=\tiny, inner sep=0pt] at (449.3,373.4) {\textarabic{٣}};
\node[font=\tiny, inner sep=0pt] at (467.6,373.4) {\EN{23/07}};
\node[font=\tiny, inner sep=0pt] at (487.1,373.4) {\EN{25/07}};
\node[font=\tiny, inner sep=0pt] at (507.3,373.4) {\EN{AI}};
\node[font=\tiny, inner sep=0pt] at (304.8,363.2) {\textarabic{١.٢.٧}};
\node[font=\scriptsize, inner sep=0pt] at (378.3,363.2) {\textarabic{النطاق وأصحاب المصلحة والقيود}};
\node[font=\tiny, inner sep=0pt] at (449.3,363.2) {\textarabic{٤}};
\node[font=\tiny, inner sep=0pt] at (467.6,363.2) {\EN{25/07}};
\node[font=\tiny, inner sep=0pt] at (487.1,363.2) {\EN{28/07}};
\node[font=\tiny, inner sep=0pt] at (507.3,363.2) {\EN{TM}};
\node[font=\tiny, inner sep=0pt] at (304.8,352.9) {\textarabic{١.٢.٨}};
\node[font=\scriptsize, inner sep=0pt] at (378.3,352.9) {\textarabic{مراجعة الفصل الأول واعتماده}};
\node[font=\tiny, inner sep=0pt] at (449.3,352.9) {\textarabic{٣}};
\node[font=\tiny, inner sep=0pt] at (467.6,352.9) {\EN{28/07}};
\node[font=\tiny, inner sep=0pt] at (487.1,352.9) {\EN{30/07}};
\node[font=\tiny, inner sep=0pt] at (507.3,352.9) {\EN{TM}};
\node[font=\tiny, inner sep=0pt] at (304.8,342.7) {\textarabic{١.٣}};
\node[font=\tiny\bfseries, inner sep=0pt] at (378.3,342.7) {\textbf{\textarabic{حزمة الفصل الثاني — مراجعة الأدبيات}}};
\node[font=\tiny, inner sep=0pt] at (449.3,342.7) {\textarabic{٢١}};
\node[font=\tiny, inner sep=0pt] at (467.6,342.7) {\EN{31/07}};
\node[font=\tiny, inner sep=0pt] at (487.1,342.7) {\EN{20/08}};
\node[font=\tiny, inner sep=0pt] at (507.3,342.7) {\textarabic{الفريق}};
\node[font=\tiny, inner sep=0pt] at (304.8,332.5) {\textarabic{١.٣.١}};
\node[font=\scriptsize, inner sep=0pt] at (378.3,332.5) {\textarabic{خطة البحث الأدبي وقائمة المراجع}};
\node[font=\tiny, inner sep=0pt] at (449.3,332.5) {\textarabic{٣}};
\node[font=\tiny, inner sep=0pt] at (467.6,332.5) {\EN{31/07}};
\node[font=\tiny, inner sep=0pt] at (487.1,332.5) {\EN{02/08}};
\node[font=\tiny, inner sep=0pt] at (507.3,332.5) {\EN{AI}};
\node[font=\tiny, inner sep=0pt] at (304.8,322.2) {\textarabic{١.٣.٢}};
\node[font=\scriptsize, inner sep=0pt] at (378.3,322.2) {\textarabic{النماذج اللغوية الكبيرة ودعم العربية}};
\node[font=\tiny, inner sep=0pt] at (449.3,322.2) {\textarabic{٣}};
\node[font=\tiny, inner sep=0pt] at (467.6,322.2) {\EN{03/08}};
\node[font=\tiny, inner sep=0pt] at (487.1,322.2) {\EN{05/08}};
\node[font=\tiny, inner sep=0pt] at (507.3,322.2) {\EN{AI}};
\node[font=\tiny, inner sep=0pt] at (304.8,312.0) {\textarabic{١.٣.٣}};
\node[font=\scriptsize, inner sep=0pt] at (378.3,312.0) {\textarabic{دراسة تقنيات التحويل إلى استعلامات}};
\node[font=\tiny, inner sep=0pt] at (449.3,312.0) {\textarabic{٤}};
\node[font=\tiny, inner sep=0pt] at (467.6,312.0) {\EN{05/08}};
\node[font=\tiny, inner sep=0pt] at (487.1,312.0) {\EN{08/08}};
\node[font=\tiny, inner sep=0pt] at (507.3,312.0) {\EN{BE}};
\node[font=\tiny, inner sep=0pt] at (304.8,301.8) {\textarabic{١.٣.٤}};
\node[font=\scriptsize, inner sep=0pt] at (378.3,301.8) {\textarabic{معمارية الاسترجاع المعزز وقواعد المتجهات}};
\node[font=\tiny, inner sep=0pt] at (449.3,301.8) {\textarabic{٣}};
\node[font=\tiny, inner sep=0pt] at (467.6,301.8) {\EN{08/08}};
\node[font=\tiny, inner sep=0pt] at (487.1,301.8) {\EN{10/08}};
\node[font=\tiny, inner sep=0pt] at (507.3,301.8) {\EN{AI}};
\node[font=\tiny, inner sep=0pt] at (304.8,291.6) {\textarabic{١.٣.٥}};
\node[font=\scriptsize, inner sep=0pt] at (378.3,291.6) {\textarabic{الأنظمة متعددة الوكلاء وأطر التنسيق}};
\node[font=\tiny, inner sep=0pt] at (449.3,291.6) {\textarabic{٣}};
\node[font=\tiny, inner sep=0pt] at (467.6,291.6) {\EN{10/08}};
\node[font=\tiny, inner sep=0pt] at (487.1,291.6) {\EN{12/08}};
\node[font=\tiny, inner sep=0pt] at (507.3,291.6) {\EN{AI}};
\node[font=\tiny, inner sep=0pt] at (304.8,281.3) {\textarabic{١.٣.٦}};
\node[font=\scriptsize, inner sep=0pt] at (378.3,281.3) {\textarabic{تحديات معالجة اللغة العربية}};
\node[font=\tiny, inner sep=0pt] at (449.3,281.3) {\textarabic{٣}};
\node[font=\tiny, inner sep=0pt] at (467.6,281.3) {\EN{12/08}};
\node[font=\tiny, inner sep=0pt] at (487.1,281.3) {\EN{14/08}};
\node[font=\tiny, inner sep=0pt] at (507.3,281.3) {\EN{AI}};
\node[font=\tiny, inner sep=0pt] at (304.8,271.1) {\textarabic{١.٣.٧}};
\node[font=\scriptsize, inner sep=0pt] at (378.3,271.1) {\textarabic{الأنظمة السابقة ومساعدات التعليم العالي}};
\node[font=\tiny, inner sep=0pt] at (449.3,271.1) {\textarabic{٤}};
\node[font=\tiny, inner sep=0pt] at (467.6,271.1) {\EN{13/08}};
\node[font=\tiny, inner sep=0pt] at (487.1,271.1) {\EN{16/08}};
\node[font=\tiny, inner sep=0pt] at (507.3,271.1) {\EN{FE}};
\node[font=\tiny, inner sep=0pt] at (304.8,260.9) {\textarabic{١.٣.٨}};
\node[font=\scriptsize, inner sep=0pt] at (378.3,260.9) {\textarabic{الجدول المقارن التحليلي}};
\node[font=\tiny, inner sep=0pt] at (449.3,260.9) {\textarabic{٢}};
\node[font=\tiny, inner sep=0pt] at (467.6,260.9) {\EN{16/08}};
\node[font=\tiny, inner sep=0pt] at (487.1,260.9) {\EN{17/08}};
\node[font=\tiny, inner sep=0pt] at (507.3,260.9) {\EN{TM}};
\node[font=\tiny, inner sep=0pt] at (304.8,250.6) {\textarabic{١.٣.٩}};
\node[font=\scriptsize, inner sep=0pt] at (378.3,250.6) {\textarabic{الفجوة البحثية والتطبيقية}};
\node[font=\tiny, inner sep=0pt] at (449.3,250.6) {\textarabic{٢}};
\node[font=\tiny, inner sep=0pt] at (467.6,250.6) {\EN{17/08}};
\node[font=\tiny, inner sep=0pt] at (487.1,250.6) {\EN{18/08}};
\node[font=\tiny, inner sep=0pt] at (507.3,250.6) {\EN{AI}};
\node[font=\tiny, inner sep=0pt] at (304.8,240.4) {\textarabic{١.٣.١٠}};
\node[font=\scriptsize, inner sep=0pt] at (378.3,240.4) {\textarabic{معايير تقييم النظام المسبقة}};
\node[font=\tiny, inner sep=0pt] at (449.3,240.4) {\textarabic{٢}};
\node[font=\tiny, inner sep=0pt] at (467.6,240.4) {\EN{18/08}};
\node[font=\tiny, inner sep=0pt] at (487.1,240.4) {\EN{19/08}};
\node[font=\tiny, inner sep=0pt] at (507.3,240.4) {\EN{AI}};
\node[font=\tiny, inner sep=0pt] at (304.8,230.2) {\textarabic{١.٣.١١}};
\node[font=\scriptsize, inner sep=0pt] at (378.3,230.2) {\textarabic{مراجعة الفصل الثاني واعتماده}};
\node[font=\tiny, inner sep=0pt] at (449.3,230.2) {\textarabic{٢}};
\node[font=\tiny, inner sep=0pt] at (467.6,230.2) {\EN{19/08}};
\node[font=\tiny, inner sep=0pt] at (487.1,230.2) {\EN{20/08}};
\node[font=\tiny, inner sep=0pt] at (507.3,230.2) {\EN{TM}};
\node[font=\tiny, inner sep=0pt] at (304.8,219.9) {\textarabic{١.٤}};
\node[font=\tiny\bfseries, inner sep=0pt] at (378.3,219.9) {\textbf{\textarabic{حزمة الفصل الثالث — المنهجية والتحليل}}};
\node[font=\tiny, inner sep=0pt] at (449.3,219.9) {\textarabic{٢٤}};
\node[font=\tiny, inner sep=0pt] at (467.6,219.9) {\EN{21/08}};
\node[font=\tiny, inner sep=0pt] at (487.1,219.9) {\EN{13/09}};
\node[font=\tiny, inner sep=0pt] at (507.3,219.9) {\textarabic{الفريق}};
\node[font=\tiny, inner sep=0pt] at (304.8,209.7) {\textarabic{١.٤.١}};
\node[font=\scriptsize, inner sep=0pt] at (378.3,209.7) {\textarabic{اختيار المنهجية الرشيقة وتبريرها}};
\node[font=\tiny, inner sep=0pt] at (449.3,209.7) {\textarabic{٣}};
\node[font=\tiny, inner sep=0pt] at (467.6,209.7) {\EN{21/08}};
\node[font=\tiny, inner sep=0pt] at (487.1,209.7) {\EN{23/08}};
\node[font=\tiny, inner sep=0pt] at (507.3,209.7) {\EN{TM}};
\node[font=\tiny, inner sep=0pt] at (304.8,199.5) {\textarabic{١.٤.٢}};
\node[font=\scriptsize, inner sep=0pt] at (378.3,199.5) {\textarabic{خارطة الدورات التنفيذية الأربع}};
\node[font=\tiny, inner sep=0pt] at (449.3,199.5) {\textarabic{٢}};
\node[font=\tiny, inner sep=0pt] at (467.6,199.5) {\EN{23/08}};
\node[font=\tiny, inner sep=0pt] at (487.1,199.5) {\EN{24/08}};
\node[font=\tiny, inner sep=0pt] at (507.3,199.5) {\EN{TM}};
\node[font=\tiny, inner sep=0pt] at (304.8,189.3) {\textarabic{١.٤.٣}};
\node[font=\scriptsize, inner sep=0pt] at (378.3,189.3) {\textarabic{دراسة الجدوى الاقتصادية}};
\node[font=\tiny, inner sep=0pt] at (449.3,189.3) {\textarabic{٤}};
\node[font=\tiny, inner sep=0pt] at (467.6,189.3) {\EN{24/08}};
\node[font=\tiny, inner sep=0pt] at (487.1,189.3) {\EN{27/08}};
\node[font=\tiny, inner sep=0pt] at (507.3,189.3) {\EN{BE}};
\node[font=\tiny, inner sep=0pt] at (304.8,179.0) {\textarabic{١.٤.٤}};
\node[font=\scriptsize, inner sep=0pt] at (378.3,179.0) {\textarabic{الجدوى التقنية والتشغيلية والقانونية}};
\node[font=\tiny, inner sep=0pt] at (449.3,179.0) {\textarabic{٤}};
\node[font=\tiny, inner sep=0pt] at (467.6,179.0) {\EN{27/08}};
\node[font=\tiny, inner sep=0pt] at (487.1,179.0) {\EN{30/08}};
\node[font=\tiny, inner sep=0pt] at (507.3,179.0) {\EN{BE}};
\node[font=\tiny, inner sep=0pt] at (304.8,168.8) {\textarabic{١.٤.٥}};
\node[font=\scriptsize, inner sep=0pt] at (378.3,168.8) {\textarabic{جمع المتطلبات من أصحاب المصلحة}};
\node[font=\tiny, inner sep=0pt] at (449.3,168.8) {\textarabic{٣}};
\node[font=\tiny, inner sep=0pt] at (467.6,168.8) {\EN{30/08}};
\node[font=\tiny, inner sep=0pt] at (487.1,168.8) {\EN{01/09}};
\node[font=\tiny, inner sep=0pt] at (507.3,168.8) {\EN{TM}};
\node[font=\tiny, inner sep=0pt] at (304.8,158.6) {\textarabic{١.٤.٦}};
\node[font=\scriptsize, inner sep=0pt] at (378.3,158.6) {\textarabic{توثيق المتطلبات الوظيفية الثلاثين}};
\node[font=\tiny, inner sep=0pt] at (449.3,158.6) {\textarabic{٣}};
\node[font=\tiny, inner sep=0pt] at (467.6,158.6) {\EN{01/09}};
\node[font=\tiny, inner sep=0pt] at (487.1,158.6) {\EN{03/09}};
\node[font=\tiny, inner sep=0pt] at (507.3,158.6) {\EN{BE}};
\node[font=\tiny, inner sep=0pt] at (304.8,148.3) {\textarabic{١.٤.٧}};
\node[font=\scriptsize, inner sep=0pt] at (378.3,148.3) {\textarabic{توثيق المتطلبات غير الوظيفية العشرين}};
\node[font=\tiny, inner sep=0pt] at (449.3,148.3) {\textarabic{٣}};
\node[font=\tiny, inner sep=0pt] at (467.6,148.3) {\EN{03/09}};
\node[font=\tiny, inner sep=0pt] at (487.1,148.3) {\EN{05/09}};
\node[font=\tiny, inner sep=0pt] at (507.3,148.3) {\EN{BE}};
\node[font=\tiny, inner sep=0pt] at (304.8,138.1) {\textarabic{١.٤.٨}};
\node[font=\scriptsize, inner sep=0pt] at (378.3,138.1) {\textarabic{بيئة التشغيل والقيود والافتراضات}};
\node[font=\tiny, inner sep=0pt] at (449.3,138.1) {\textarabic{٢}};
\node[font=\tiny, inner sep=0pt] at (467.6,138.1) {\EN{05/09}};
\node[font=\tiny, inner sep=0pt] at (487.1,138.1) {\EN{06/09}};
\node[font=\tiny, inner sep=0pt] at (507.3,138.1) {\EN{BE}};
\node[font=\tiny, inner sep=0pt] at (304.8,127.9) {\textarabic{١.٤.٩}};
\node[font=\scriptsize, inner sep=0pt] at (378.3,127.9) {\textarabic{الفاعلون ومخطط حالات الاستخدام}};
\node[font=\tiny, inner sep=0pt] at (449.3,127.9) {\textarabic{٣}};
\node[font=\tiny, inner sep=0pt] at (467.6,127.9) {\EN{06/09}};
\node[font=\tiny, inner sep=0pt] at (487.1,127.9) {\EN{08/09}};
\node[font=\tiny, inner sep=0pt] at (507.3,127.9) {\EN{FE}};
\node[font=\tiny, inner sep=0pt] at (304.8,117.6) {\textarabic{١.٤.١٠}};
\node[font=\scriptsize, inner sep=0pt] at (378.3,117.6) {\textarabic{أوصاف حالات الاستخدام الثماني}};
\node[font=\tiny, inner sep=0pt] at (449.3,117.6) {\textarabic{٣}};
\node[font=\tiny, inner sep=0pt] at (467.6,117.6) {\EN{08/09}};
\node[font=\tiny, inner sep=0pt] at (487.1,117.6) {\EN{10/09}};
\node[font=\tiny, inner sep=0pt] at (507.3,117.6) {\EN{FE}};
\node[font=\tiny, inner sep=0pt] at (304.8,107.4) {\textarabic{١.٤.١١}};
\node[font=\scriptsize, inner sep=0pt] at (378.3,107.4) {\textarabic{مخطط التسلسل الشامل}};
\node[font=\tiny, inner sep=0pt] at (449.3,107.4) {\textarabic{٢}};
\node[font=\tiny, inner sep=0pt] at (467.6,107.4) {\EN{10/09}};
\node[font=\tiny, inner sep=0pt] at (487.1,107.4) {\EN{11/09}};
\node[font=\tiny, inner sep=0pt] at (507.3,107.4) {\EN{BE}};
\node[font=\tiny, inner sep=0pt] at (304.8,97.2) {\textarabic{١.٤.١٢}};
\node[font=\scriptsize, inner sep=0pt] at (378.3,97.2) {\textarabic{منهجية تحليل المخاطر الكمية}};
\node[font=\tiny, inner sep=0pt] at (449.3,97.2) {\textarabic{٢}};
\node[font=\tiny, inner sep=0pt] at (467.6,97.2) {\EN{11/09}};
\node[font=\tiny, inner sep=0pt] at (487.1,97.2) {\EN{12/09}};
\node[font=\tiny, inner sep=0pt] at (507.3,97.2) {\EN{TM}};
\node[font=\tiny, inner sep=0pt] at (304.8,87.0) {\textarabic{١.٤.١٣}};
\node[font=\scriptsize, inner sep=0pt] at (378.3,87.0) {\textarabic{سجل المخاطر ومصفوفتها}};
\node[font=\tiny, inner sep=0pt] at (449.3,87.0) {\textarabic{٢}};
\node[font=\tiny, inner sep=0pt] at (467.6,87.0) {\EN{12/09}};
\node[font=\tiny, inner sep=0pt] at (487.1,87.0) {\EN{13/09}};
\node[font=\tiny, inner sep=0pt] at (507.3,87.0) {\EN{TM}};
\node[font=\tiny, inner sep=0pt] at (304.8,76.7) {\textarabic{١.٤.١٤}};
\node[font=\scriptsize, inner sep=0pt] at (378.3,76.7) {\textarabic{خطط الاستجابة واعتماد الفصل}};
\node[font=\tiny, inner sep=0pt] at (449.3,76.7) {\textarabic{٢}};
\node[font=\tiny, inner sep=0pt] at (467.6,76.7) {\EN{12/09}};
\node[font=\tiny, inner sep=0pt] at (487.1,76.7) {\EN{13/09}};
\node[font=\tiny, inner sep=0pt] at (507.3,76.7) {\EN{TM}};
\node[font=\tiny, inner sep=0pt] at (304.8,66.5) {\textarabic{١.٥}};
\node[font=\tiny\bfseries, inner sep=0pt] at (378.3,66.5) {\textbf{\textarabic{حزمة الإصدار الأول والتسليم}}};
\node[font=\tiny, inner sep=0pt] at (449.3,66.5) {\textarabic{٢٧}};
\node[font=\tiny, inner sep=0pt] at (467.6,66.5) {\EN{01/09}};
\node[font=\tiny, inner sep=0pt] at (487.1,66.5) {\EN{27/09}};
\node[font=\tiny, inner sep=0pt] at (507.3,66.5) {\textarabic{الفريق}};
\node[font=\tiny, inner sep=0pt] at (304.8,56.3) {\textarabic{١.٥.١}};
\node[font=\scriptsize, inner sep=0pt] at (378.3,56.3) {\textarabic{التدقيق اللغوي والفني المتقاطع}};
\node[font=\tiny, inner sep=0pt] at (449.3,56.3) {\textarabic{٦}};
\node[font=\tiny, inner sep=0pt] at (467.6,56.3) {\EN{01/09}};
\node[font=\tiny, inner sep=0pt] at (487.1,56.3) {\EN{06/09}};
\node[font=\tiny, inner sep=0pt] at (507.3,56.3) {\EN{TM}};
\node[font=\tiny, inner sep=0pt] at (304.8,46.0) {\textarabic{١.٥.٢}};
\node[font=\scriptsize, inner sep=0pt] at (378.3,46.0) {\textarabic{ضبط التنسيق وبناء النسخة النهائية}};
\node[font=\tiny, inner sep=0pt] at (449.3,46.0) {\textarabic{٦}};
\node[font=\tiny, inner sep=0pt] at (467.6,46.0) {\EN{05/09}};
\node[font=\tiny, inner sep=0pt] at (487.1,46.0) {\EN{10/09}};
\node[font=\tiny, inner sep=0pt] at (507.3,46.0) {\EN{TM}};
\node[font=\tiny, inner sep=0pt] at (304.8,35.8) {\textarabic{١.٥.٣}};
\node[font=\scriptsize, inner sep=0pt] at (378.3,35.8) {\textarabic{تسليم نسختي الوثيقة إلى الكلية}};
\node[font=\tiny, inner sep=0pt] at (449.3,35.8) {\textarabic{٣}};
\node[font=\tiny, inner sep=0pt] at (467.6,35.8) {\EN{11/09}};
\node[font=\tiny, inner sep=0pt] at (487.1,35.8) {\EN{13/09}};
\node[font=\tiny, inner sep=0pt] at (507.3,35.8) {\EN{TM}};
\node[font=\tiny, inner sep=0pt] at (304.8,25.6) {\textarabic{١.٥.٤}};
\node[font=\scriptsize, inner sep=0pt] at (378.3,25.6) {\textarabic{إعداد عرض المناقشة الأولى}};
\node[font=\tiny, inner sep=0pt] at (449.3,25.6) {\textarabic{٣}};
\node[font=\tiny, inner sep=0pt] at (467.6,25.6) {\EN{14/09}};
\node[font=\tiny, inner sep=0pt] at (487.1,25.6) {\EN{16/09}};
\node[font=\tiny, inner sep=0pt] at (507.3,25.6) {\EN{FE}};
\node[font=\tiny, inner sep=0pt] at (304.8,15.3) {\textarabic{١.٥.٥}};
\node[font=\scriptsize, inner sep=0pt] at (378.3,15.3) {\textarabic{المحاكاة الداخلية للمناقشة}};
\node[font=\tiny, inner sep=0pt] at (449.3,15.3) {\textarabic{٣}};
\node[font=\tiny, inner sep=0pt] at (467.6,15.3) {\EN{17/09}};
\node[font=\tiny, inner sep=0pt] at (487.1,15.3) {\EN{19/09}};
\node[font=\tiny, inner sep=0pt] at (507.3,15.3) {\EN{TM}};
\node[font=\tiny, inner sep=0pt] at (304.8,5.1) {\textarabic{١.٥.٦}};
\node[font=\scriptsize, inner sep=0pt] at (378.3,5.1) {\textarabic{المناقشة النهائية — بوابة الجودة الأولى}};
\node[font=\tiny, inner sep=0pt] at (449.3,5.1) {\textarabic{٨}};
\node[font=\tiny, inner sep=0pt] at (467.6,5.1) {\EN{20/09}};
\node[font=\tiny, inner sep=0pt] at (487.1,5.1) {\EN{27/09}};
\node[font=\tiny, inner sep=0pt] at (507.3,5.1) {\EN{TM}};
\end{tikzpicture}%
\end{LTR}%
\end{lrbox}%
\noindent\makebox[\textwidth][c]{\usebox\gCBA}%
\begin{center}
\scriptsize
\renewcommand{\arraystretch}{1.0}
\begin{tabular}{c@{\hspace{3pt}}l@{\hspace{7pt}}c@{\hspace{3pt}}l@{\hspace{7pt}}c@{\hspace{3pt}}l}
  \tikz\fill[ownerAI!55,draw=ownerAI] (0,0) rectangle (0.28,0.14); & الذكاء الاصطناعي (\EN{AI}) &
  \tikz\fill[ownerBE!55,draw=ownerBE] (0,0) rectangle (0.28,0.14); & الخلفية وقواعد البيانات (\EN{BE}) &
  \tikz\fill[ownerFE!45,draw=ownerFE] (0,0) rectangle (0.28,0.14); & الواجهة وتجربة المستخدم (\EN{FE}) \\[3pt]
  \tikz\fill[ownerTM!50,draw=ownerTM] (0,0) rectangle (0.28,0.14); & الفريق كاملاً (\EN{TM}) &
  \tikz\fill[kaudark!90,draw=kaudark!90] (0,0.05) rectangle (0.4,0.11); & شريط الحزمة (الملخص) &
  \tikz{\fill[kauaccent,draw=kauaccent!50!black] (0.13,0.26) -- (0.26,0.13) -- (0.13,0) -- (0,0.13) -- cycle;} & معلم رسمي بتاريخه الرسمي \\[3pt]
  \tikz\fill[kauaccent!35,draw=kauaccent!70!black] (0,0) rectangle (0.28,0.14); & إعداد العرض والمحاكاة &
  \tikz\fill[kauaccent!75,draw=kauaccent!80!black] (0,0) rectangle (0.28,0.14); & أسبوع المناقشة النهائية &
  \tikz{\draw[-stealth,thick,kauprimary] (0,0.07) -- (0.4,0.07);} & سهم تبعية (المسار الحرج) \\[3pt]
  \tikz\fill[gray!40] (0,0.02) rectangle (0.3,0.13); & ظل عطلة نهاية الأسبوع &
  \tikz{\draw[gray!45] (0.02,0) -- (0.02,0.16);} & بداية الأسبوع (الأحد) &
  \tikz\fill[kauprimary!22,draw=kauprimary!55] (0,0) rectangle (0.3,0.14); & ترويسة لوحة البيانات \\
\end{tabular}
\end{center}
  \caption{مخطط جانت التنفيذي للمرحلة الأولى — مشروع التخرج 1 (يوليو -- سبتمبر 2026): لوحة موحَّدة لـ44 مهمة بأسلوب برامج إدارة المشاريع مع المعالم \EN{M1}--\EN{M6} وأسهم المسار الحرج، ويتقدم الزمن فيها من اليمين الملاصق للوحة البيانات نحو اليسار.}
  \label{fig:ganttP1}
\end{figure}


\begin{figure}[H]
  \centering
\definecolor{ownerAI}{HTML}{0E7490}%
\definecolor{ownerBE}{HTML}{2E75B6}%
\definecolor{ownerFE}{HTML}{6D28D9}%
\definecolor{ownerTM}{HTML}{64748B}%
\newsavebox\gCBB
\begin{lrbox}{\gCBB}%
\begin{LTR}%
\begin{tikzpicture}[x=1pt, y=1pt]
\fill[kauprimary] (0,482.2) rectangle (517.8,495.0);
\fill[kauprimary!14] (0,470.6) rectangle (293.8,482.2);
\fill[kauprimary!7] (0,460.4) rectangle (293.8,470.6);
\fill[kauprimary!22] (293.8,460.4) rectangle (517.8,482.2);
\fill[kauprimary!10] (0,450.1) rectangle (517.8,460.4);
\fill[kauprimary!10] (0,388.7) rectangle (517.8,399.0);
\fill[kauprimary!10] (0,306.9) rectangle (517.8,317.1);
\fill[kauprimary!10] (0,204.6) rectangle (517.8,214.8);
\fill[kauprimary!10] (0,122.8) rectangle (517.8,133.0);
\fill[kauprimary!10] (0,71.6) rectangle (517.8,81.8);
\fill[gray!38] (270.1,470.6) rectangle (276.0,0.0);
\fill[gray!38] (249.3,470.6) rectangle (255.3,0.0);
\fill[gray!38] (228.5,470.6) rectangle (234.5,0.0);
\fill[gray!38] (207.8,470.6) rectangle (213.7,0.0);
\fill[gray!38] (187.0,470.6) rectangle (192.9,0.0);
\fill[gray!38] (166.2,470.6) rectangle (172.1,0.0);
\fill[gray!38] (145.4,470.6) rectangle (151.4,0.0);
\fill[gray!38] (124.7,470.6) rectangle (130.6,0.0);
\fill[gray!38] (103.9,470.6) rectangle (109.8,0.0);
\fill[gray!38] (83.1,470.6) rectangle (89.0,0.0);
\fill[gray!38] (62.3,470.6) rectangle (68.3,0.0);
\fill[gray!38] (41.6,470.6) rectangle (47.5,0.0);
\fill[gray!38] (20.8,470.6) rectangle (26.7,0.0);
\fill[gray!38] (0.0,470.6) rectangle (5.9,0.0);
\draw[gray!55, line width=0.45pt] (249.3,0) -- (249.3,482.2);
\draw[gray!55, line width=0.45pt] (160.3,0) -- (160.3,482.2);
\draw[gray!55, line width=0.45pt] (68.3,0) -- (68.3,482.2);
\draw[gray!45, line width=0.3pt] (290.9,0) -- (290.9,470.6);
\draw[gray!45, line width=0.3pt] (270.1,0) -- (270.1,470.6);
\draw[gray!45, line width=0.3pt] (249.3,0) -- (249.3,470.6);
\draw[gray!45, line width=0.3pt] (228.5,0) -- (228.5,470.6);
\draw[gray!45, line width=0.3pt] (207.8,0) -- (207.8,470.6);
\draw[gray!45, line width=0.3pt] (187.0,0) -- (187.0,470.6);
\draw[gray!45, line width=0.3pt] (166.2,0) -- (166.2,470.6);
\draw[gray!45, line width=0.3pt] (145.4,0) -- (145.4,470.6);
\draw[gray!45, line width=0.3pt] (124.7,0) -- (124.7,470.6);
\draw[gray!45, line width=0.3pt] (103.9,0) -- (103.9,470.6);
\draw[gray!45, line width=0.3pt] (83.1,0) -- (83.1,470.6);
\draw[gray!45, line width=0.3pt] (62.3,0) -- (62.3,470.6);
\draw[gray!45, line width=0.3pt] (41.6,0) -- (41.6,470.6);
\draw[gray!45, line width=0.3pt] (20.8,0) -- (20.8,470.6);
\fill[kaudark!90] (276.0,453.8) rectangle (293.8,456.1);
\filldraw[fill=ownerTM!50, draw=ownerTM, line width=0.4pt] (287.9,441.9) rectangle (293.8,448.1);
\filldraw[fill=ownerBE!55, draw=ownerBE, line width=0.4pt] (282.0,431.7) rectangle (293.8,437.8);
\filldraw[fill=ownerBE!55, draw=ownerBE, line width=0.4pt] (279.0,421.5) rectangle (287.9,427.6);
\filldraw[fill=ownerAI!55, draw=ownerAI, line width=0.4pt] (276.0,411.2) rectangle (284.9,417.4);
\filldraw[fill=ownerBE!55, draw=ownerBE, line width=0.4pt] (276.0,401.0) rectangle (282.0,407.2);
\fill[kaudark!90] (207.8,392.4) rectangle (276.0,394.7);
\filldraw[fill=ownerAI!55, draw=ownerAI, line width=0.4pt] (261.2,380.6) rectangle (276.0,386.7);
\filldraw[fill=ownerBE!55, draw=ownerBE, line width=0.4pt] (249.3,370.3) rectangle (264.2,376.5);
\filldraw[fill=ownerFE!45, draw=ownerFE, line width=0.4pt] (237.4,360.1) rectangle (252.3,366.2);
\filldraw[fill=ownerTM!50, draw=ownerTM, line width=0.4pt] (228.5,349.9) rectangle (243.4,356.0);
\filldraw[fill=ownerBE!55, draw=ownerBE, line width=0.4pt] (222.6,339.6) rectangle (234.5,345.8);
\filldraw[fill=ownerBE!55, draw=ownerBE, line width=0.4pt] (213.7,329.4) rectangle (225.6,335.5);
\filldraw[fill=ownerTM!50, draw=ownerTM, line width=0.4pt] (207.8,319.2) rectangle (216.7,325.3);
\fill[kaudark!90] (118.7,310.6) rectangle (207.8,312.8);
\filldraw[fill=ownerAI!55, draw=ownerAI, line width=0.4pt] (184.0,298.7) rectangle (207.8,304.9);
\filldraw[fill=ownerAI!55, draw=ownerAI, line width=0.4pt] (160.3,288.5) rectangle (184.0,294.6);
\filldraw[fill=ownerFE!45, draw=ownerFE, line width=0.4pt] (136.5,278.3) rectangle (160.3,284.4);
\filldraw[fill=ownerBE!55, draw=ownerBE, line width=0.4pt] (154.3,268.0) rectangle (207.8,274.2);
\filldraw[fill=ownerBE!55, draw=ownerBE, line width=0.4pt] (136.5,257.8) rectangle (181.1,263.9);
\filldraw[fill=ownerBE!55, draw=ownerBE, line width=0.4pt] (118.7,247.6) rectangle (160.3,253.7);
\filldraw[fill=ownerBE!55, draw=ownerBE, line width=0.4pt] (118.7,237.3) rectangle (207.8,243.5);
\filldraw[fill=ownerBE!55, draw=ownerBE, line width=0.4pt] (118.7,227.1) rectangle (136.5,233.2);
\filldraw[fill=ownerTM!50, draw=ownerTM, line width=0.4pt] (118.7,216.9) rectangle (127.6,223.0);
\fill[kaudark!90] (32.6,208.3) rectangle (118.7,210.5);
\filldraw[fill=ownerAI!55, draw=ownerAI, line width=0.4pt] (97.9,196.4) rectangle (118.7,202.6);
\filldraw[fill=ownerAI!55, draw=ownerAI, line width=0.4pt] (77.2,186.2) rectangle (97.9,192.3);
\filldraw[fill=ownerBE!55, draw=ownerBE, line width=0.4pt] (59.4,176.0) rectangle (77.2,182.1);
\filldraw[fill=ownerBE!55, draw=ownerBE, line width=0.4pt] (53.4,165.7) rectangle (74.2,171.9);
\filldraw[fill=ownerFE!45, draw=ownerFE, line width=0.4pt] (44.5,155.5) rectangle (59.4,161.6);
\filldraw[fill=ownerTM!50, draw=ownerTM, line width=0.4pt] (38.6,145.3) rectangle (56.4,151.4);
\filldraw[fill=ownerTM!50, draw=ownerTM, line width=0.4pt] (32.6,135.0) rectangle (41.6,141.2);
\fill[kaudark!90] (32.6,126.4) rectangle (68.3,128.7);
\filldraw[fill=ownerAI!55, draw=ownerAI, line width=0.4pt] (56.4,114.6) rectangle (68.3,120.7);
\filldraw[fill=ownerTM!50, draw=ownerTM, line width=0.4pt] (50.5,104.3) rectangle (59.4,110.5);
\filldraw[fill=ownerTM!50, draw=ownerTM, line width=0.4pt] (44.5,94.1) rectangle (53.4,100.3);
\filldraw[fill=ownerTM!50, draw=ownerTM, line width=0.4pt] (32.6,83.9) rectangle (41.6,90.0);
\fill[kaudark!90] (0.0,75.3) rectangle (68.3,77.5);
\filldraw[fill=ownerTM!50, draw=ownerTM, line width=0.4pt] (53.4,63.4) rectangle (68.3,69.6);
\filldraw[fill=ownerFE!45, draw=ownerFE, line width=0.4pt] (47.5,53.2) rectangle (59.4,59.3);
\filldraw[fill=ownerFE!45, draw=ownerFE, line width=0.4pt] (41.6,43.0) rectangle (53.4,49.1);
\filldraw[fill=ownerTM!50, draw=ownerTM, line width=0.4pt] (41.6,32.7) rectangle (47.5,38.9);
\filldraw[fill=kauaccent!35, draw=kauaccent!70!black, line width=0.4pt] (23.7,22.5) rectangle (41.6,28.6);
\filldraw[fill=kauaccent!75, draw=kauaccent!80!black, line width=0.4pt] (0.0,12.3) rectangle (23.7,18.4);
\filldraw[fill=kauaccent!75, draw=kauaccent!80!black, line width=0.4pt] (0.0,2.0) rectangle (11.9,8.2);
\node[diamond, fill=kauaccent, draw=kauaccent!60, inner sep=2.6pt, line width=0.4pt] at (283.4,445.0) {};
\node[anchor=east, font=\tiny\bfseries, inner sep=0pt] at (279.2,445.0) {\EN{M8 20/10}};
\node[diamond, fill=kauaccent, draw=kauaccent!60, inner sep=2.6pt, line width=0.4pt] at (209.2,322.2) {};
\node[anchor=east, font=\tiny\bfseries, inner sep=0pt] at (205.0,322.2) {\EN{M9 14/11}};
\node[diamond, fill=kauaccent, draw=kauaccent!60, inner sep=2.6pt, line width=0.4pt] at (120.2,219.9) {};
\node[anchor=east, font=\tiny\bfseries, inner sep=0pt] at (116.0,219.9) {\EN{M10 14/12}};
\node[diamond, fill=kauaccent, draw=kauaccent!60, inner sep=2.6pt, line width=0.4pt] at (34.1,138.1) {};
\node[anchor=south, font=\tiny\bfseries, inner sep=0pt] at (34.1,140.6) {\EN{M12 12/01}};
\node[diamond, fill=kauaccent, draw=kauaccent!60, inner sep=2.6pt, line width=0.4pt] at (43.0,46.0) {};
\node[anchor=south, font=\tiny\bfseries, inner sep=0pt] at (43.0,48.5) {\EN{M11 09/01}};
\node[diamond, fill=kauaccent, draw=kauaccent!60, inner sep=2.6pt, line width=0.4pt] at (22.3,25.6) {};
\node[anchor=south, font=\tiny\bfseries, inner sep=0pt] at (22.3,28.1) {\EN{M13 16/01}};
\draw[-stealth, kauprimary, line width=0.75pt] (207.8,393.9) -- (204.8,393.9) -- (204.8,301.8) -- (184.0,301.8);
\draw[-stealth, kauprimary, line width=0.75pt] (184.0,301.8) -- (181.0,301.8) -- (181.0,291.6) -- (160.3,291.6);
\draw[-stealth, kauprimary, line width=0.75pt] (160.3,291.6) -- (157.3,291.6) -- (157.3,281.3) -- (136.5,281.3);
\draw[-stealth, kauprimary, line width=0.75pt] (136.5,281.3) -- (133.5,281.3) -- (133.5,230.2) -- (118.7,230.2);
\draw[-stealth, kauprimary, line width=0.75pt] (118.7,230.2) -- (115.7,230.2) -- (115.7,209.7) -- (32.6,209.7);
\draw[-stealth, kauprimary, line width=0.75pt] (32.6,209.7) -- (29.6,209.7) -- (29.6,15.3) -- (0.0,15.3);
\draw[kauprimary!50, line width=0.5pt] (0,460.4) rectangle (517.8,460.4);
\draw[kauprimary!45, line width=0.45pt] (0,450.1) -- (517.8,450.1);
\draw[gray!32, line width=0.3pt] (0,439.9) -- (517.8,439.9);
\draw[gray!32, line width=0.3pt] (0,429.7) -- (517.8,429.7);
\draw[gray!32, line width=0.3pt] (0,419.4) -- (517.8,419.4);
\draw[gray!32, line width=0.3pt] (0,409.2) -- (517.8,409.2);
\draw[gray!32, line width=0.3pt] (0,399.0) -- (517.8,399.0);
\draw[kauprimary!45, line width=0.45pt] (0,388.7) -- (517.8,388.7);
\draw[gray!32, line width=0.3pt] (0,378.5) -- (517.8,378.5);
\draw[gray!32, line width=0.3pt] (0,368.3) -- (517.8,368.3);
\draw[gray!32, line width=0.3pt] (0,358.1) -- (517.8,358.1);
\draw[gray!32, line width=0.3pt] (0,347.8) -- (517.8,347.8);
\draw[gray!32, line width=0.3pt] (0,337.6) -- (517.8,337.6);
\draw[gray!32, line width=0.3pt] (0,327.4) -- (517.8,327.4);
\draw[gray!32, line width=0.3pt] (0,317.1) -- (517.8,317.1);
\draw[kauprimary!45, line width=0.45pt] (0,306.9) -- (517.8,306.9);
\draw[gray!32, line width=0.3pt] (0,296.7) -- (517.8,296.7);
\draw[gray!32, line width=0.3pt] (0,286.4) -- (517.8,286.4);
\draw[gray!32, line width=0.3pt] (0,276.2) -- (517.8,276.2);
\draw[gray!32, line width=0.3pt] (0,266.0) -- (517.8,266.0);
\draw[gray!32, line width=0.3pt] (0,255.8) -- (517.8,255.8);
\draw[gray!32, line width=0.3pt] (0,245.5) -- (517.8,245.5);
\draw[gray!32, line width=0.3pt] (0,235.3) -- (517.8,235.3);
\draw[gray!32, line width=0.3pt] (0,225.1) -- (517.8,225.1);
\draw[gray!32, line width=0.3pt] (0,214.8) -- (517.8,214.8);
\draw[kauprimary!45, line width=0.45pt] (0,204.6) -- (517.8,204.6);
\draw[gray!32, line width=0.3pt] (0,194.4) -- (517.8,194.4);
\draw[gray!32, line width=0.3pt] (0,184.1) -- (517.8,184.1);
\draw[gray!32, line width=0.3pt] (0,173.9) -- (517.8,173.9);
\draw[gray!32, line width=0.3pt] (0,163.7) -- (517.8,163.7);
\draw[gray!32, line width=0.3pt] (0,153.4) -- (517.8,153.4);
\draw[gray!32, line width=0.3pt] (0,143.2) -- (517.8,143.2);
\draw[gray!32, line width=0.3pt] (0,133.0) -- (517.8,133.0);
\draw[kauprimary!45, line width=0.45pt] (0,122.8) -- (517.8,122.8);
\draw[gray!32, line width=0.3pt] (0,112.5) -- (517.8,112.5);
\draw[gray!32, line width=0.3pt] (0,102.3) -- (517.8,102.3);
\draw[gray!32, line width=0.3pt] (0,92.1) -- (517.8,92.1);
\draw[gray!32, line width=0.3pt] (0,81.8) -- (517.8,81.8);
\draw[kauprimary!45, line width=0.45pt] (0,71.6) -- (517.8,71.6);
\draw[gray!32, line width=0.3pt] (0,61.4) -- (517.8,61.4);
\draw[gray!32, line width=0.3pt] (0,51.1) -- (517.8,51.1);
\draw[gray!32, line width=0.3pt] (0,40.9) -- (517.8,40.9);
\draw[gray!32, line width=0.3pt] (0,30.7) -- (517.8,30.7);
\draw[gray!32, line width=0.3pt] (0,20.5) -- (517.8,20.5);
\draw[gray!32, line width=0.3pt] (0,10.2) -- (517.8,10.2);
\draw[gray!32, line width=0.3pt] (0,0.0) -- (517.8,0.0);
\draw[kauprimary!70, line width=0.6pt] (0,482.2) -- (517.8,482.2);
\draw[kauprimary!35, line width=0.4pt] (0,470.6) -- (293.8,470.6);
\draw[gray!35, line width=0.3pt] (315.8,0) -- (315.8,460.4);
\draw[kauprimary!55, line width=0.4pt] (315.8,460.4) -- (315.8,482.2);
\draw[gray!35, line width=0.3pt] (440.8,0) -- (440.8,460.4);
\draw[kauprimary!55, line width=0.4pt] (440.8,460.4) -- (440.8,482.2);
\draw[gray!35, line width=0.3pt] (457.8,0) -- (457.8,460.4);
\draw[kauprimary!55, line width=0.4pt] (457.8,460.4) -- (457.8,482.2);
\draw[gray!35, line width=0.3pt] (477.3,0) -- (477.3,460.4);
\draw[kauprimary!55, line width=0.4pt] (477.3,460.4) -- (477.3,482.2);
\draw[gray!35, line width=0.3pt] (496.8,0) -- (496.8,460.4);
\draw[kauprimary!55, line width=0.4pt] (496.8,460.4) -- (496.8,482.2);
\draw[kauprimary!60, line width=0.65pt] (293.8,0) -- (293.8,482.2);
\draw[kauprimary!70, line width=0.8pt] (0,0) rectangle (517.8,495.0);
\node[font=\scriptsize\bfseries, text=white, inner sep=0pt] at (258.9,488.6) {\textarabic{مشروع التخرج 2 — المرحلتان الثانية والثالثة: التنفيذ والاختبار والتسليم النهائي}};
\node[font=\tiny\bfseries, inner sep=0pt] at (271.6,476.4) {\textarabic{أكتوبر 2026}};
\node[font=\tiny\bfseries, inner sep=0pt] at (204.8,476.4) {\textarabic{نوفمبر}};
\node[font=\tiny\bfseries, inner sep=0pt] at (114.3,476.4) {\textarabic{ديسمبر}};
\node[font=\tiny\bfseries, inner sep=0pt] at (34.1,476.4) {\textarabic{يناير 2027}};
\node[font=\tiny, inner sep=0pt] at (280.5,465.5) {\textarabic{١}};
\node[font=\tiny, inner sep=0pt] at (259.7,465.5) {\textarabic{٢}};
\node[font=\tiny, inner sep=0pt] at (238.9,465.5) {\textarabic{٣}};
\node[font=\tiny, inner sep=0pt] at (218.2,465.5) {\textarabic{٤}};
\node[font=\tiny, inner sep=0pt] at (197.4,465.5) {\textarabic{٥}};
\node[font=\tiny, inner sep=0pt] at (176.6,465.5) {\textarabic{٦}};
\node[font=\tiny, inner sep=0pt] at (155.8,465.5) {\textarabic{٧}};
\node[font=\tiny, inner sep=0pt] at (135.0,465.5) {\textarabic{٨}};
\node[font=\tiny, inner sep=0pt] at (114.3,465.5) {\textarabic{٩}};
\node[font=\tiny, inner sep=0pt] at (93.5,465.5) {\textarabic{١٠}};
\node[font=\tiny, inner sep=0pt] at (72.7,465.5) {\textarabic{١١}};
\node[font=\tiny, inner sep=0pt] at (51.9,465.5) {\textarabic{١٢}};
\node[font=\tiny, inner sep=0pt] at (31.2,465.5) {\textarabic{١٣}};
\node[font=\tiny, inner sep=0pt] at (10.4,465.5) {\textarabic{١٤}};
\node[font=\tiny\bfseries, inner sep=0pt] at (304.8,471.3) {\textarabic{الرمز}};
\node[font=\tiny\bfseries, inner sep=0pt] at (378.3,471.3) {\textarabic{المهمة}};
\node[font=\tiny\bfseries, inner sep=0pt] at (449.3,471.3) {\textarabic{الأيام}};
\node[font=\tiny\bfseries, inner sep=0pt] at (467.6,471.3) {\textarabic{البداية}};
\node[font=\tiny\bfseries, inner sep=0pt] at (487.1,471.3) {\textarabic{النهاية}};
\node[font=\tiny\bfseries, inner sep=0pt] at (507.3,471.3) {\textarabic{المسؤول}};
\node[font=\tiny, inner sep=0pt] at (304.8,455.2) {\textarabic{٢.١}};
\node[font=\tiny\bfseries, inner sep=0pt] at (378.3,455.2) {\textbf{\textarabic{حزمة التهيئة التقنية}}};
\node[font=\tiny, inner sep=0pt] at (449.3,455.2) {\textarabic{٦}};
\node[font=\tiny, inner sep=0pt] at (467.6,455.2) {\EN{17/10}};
\node[font=\tiny, inner sep=0pt] at (487.1,455.2) {\EN{22/10}};
\node[font=\tiny, inner sep=0pt] at (507.3,455.2) {\textarabic{الفريق}};
\node[font=\tiny, inner sep=0pt] at (304.8,445.0) {\textarabic{٢.١.١}};
\node[font=\scriptsize, inner sep=0pt] at (378.3,445.0) {\textarabic{استيعاب ملاحظات المناقشة الأولى}};
\node[font=\tiny, inner sep=0pt] at (449.3,445.0) {\textarabic{٢}};
\node[font=\tiny, inner sep=0pt] at (467.6,445.0) {\EN{17/10}};
\node[font=\tiny, inner sep=0pt] at (487.1,445.0) {\EN{18/10}};
\node[font=\tiny, inner sep=0pt] at (507.3,445.0) {\EN{TM}};
\node[font=\tiny, inner sep=0pt] at (304.8,434.8) {\textarabic{٢.١.٢}};
\node[font=\scriptsize, inner sep=0pt] at (378.3,434.8) {\textarabic{بيئة الحاويات الموحدة للخدمات}};
\node[font=\tiny, inner sep=0pt] at (449.3,434.8) {\textarabic{٤}};
\node[font=\tiny, inner sep=0pt] at (467.6,434.8) {\EN{17/10}};
\node[font=\tiny, inner sep=0pt] at (487.1,434.8) {\EN{20/10}};
\node[font=\tiny, inner sep=0pt] at (507.3,434.8) {\EN{BE}};
\node[font=\tiny, inner sep=0pt] at (304.8,424.5) {\textarabic{٢.١.٣}};
\node[font=\scriptsize, inner sep=0pt] at (378.3,424.5) {\textarabic{مستودع الشفرة ومعايير الجودة}};
\node[font=\tiny, inner sep=0pt] at (449.3,424.5) {\textarabic{٣}};
\node[font=\tiny, inner sep=0pt] at (467.6,424.5) {\EN{19/10}};
\node[font=\tiny, inner sep=0pt] at (487.1,424.5) {\EN{21/10}};
\node[font=\tiny, inner sep=0pt] at (507.3,424.5) {\EN{BE}};
\node[font=\tiny, inner sep=0pt] at (304.8,414.3) {\textarabic{٢.١.٤}};
\node[font=\scriptsize, inner sep=0pt] at (378.3,414.3) {\textarabic{تنزيل النماذج واختبار ملاءمتها}};
\node[font=\tiny, inner sep=0pt] at (449.3,414.3) {\textarabic{٣}};
\node[font=\tiny, inner sep=0pt] at (467.6,414.3) {\EN{20/10}};
\node[font=\tiny, inner sep=0pt] at (487.1,414.3) {\EN{22/10}};
\node[font=\tiny, inner sep=0pt] at (507.3,414.3) {\EN{AI}};
\node[font=\tiny, inner sep=0pt] at (304.8,404.1) {\textarabic{٢.١.٥}};
\node[font=\scriptsize, inner sep=0pt] at (378.3,404.1) {\textarabic{بيئة التطبيق والبيانات الاصطناعية}};
\node[font=\tiny, inner sep=0pt] at (449.3,404.1) {\textarabic{٢}};
\node[font=\tiny, inner sep=0pt] at (467.6,404.1) {\EN{21/10}};
\node[font=\tiny, inner sep=0pt] at (487.1,404.1) {\EN{22/10}};
\node[font=\tiny, inner sep=0pt] at (507.3,404.1) {\EN{BE}};
\node[font=\tiny, inner sep=0pt] at (304.8,393.9) {\textarabic{٢.٢}};
\node[font=\tiny\bfseries, inner sep=0pt] at (378.3,393.9) {\textbf{\textarabic{حزمة الفصل الرابع — تصميم النظام}}};
\node[font=\tiny, inner sep=0pt] at (449.3,393.9) {\textarabic{٢٣}};
\node[font=\tiny, inner sep=0pt] at (467.6,393.9) {\EN{23/10}};
\node[font=\tiny, inner sep=0pt] at (487.1,393.9) {\EN{14/11}};
\node[font=\tiny, inner sep=0pt] at (507.3,393.9) {\textarabic{متعدد}};
\node[font=\tiny, inner sep=0pt] at (304.8,383.6) {\textarabic{٢.٢.١}};
\node[font=\scriptsize, inner sep=0pt] at (378.3,383.6) {\textarabic{المعمارية العامة متعددة الوكلاء}};
\node[font=\tiny, inner sep=0pt] at (449.3,383.6) {\textarabic{٥}};
\node[font=\tiny, inner sep=0pt] at (467.6,383.6) {\EN{23/10}};
\node[font=\tiny, inner sep=0pt] at (487.1,383.6) {\EN{27/10}};
\node[font=\tiny, inner sep=0pt] at (507.3,383.6) {\EN{AI}};
\node[font=\tiny, inner sep=0pt] at (304.8,373.4) {\textarabic{٢.٢.٢}};
\node[font=\scriptsize, inner sep=0pt] at (378.3,373.4) {\textarabic{قواعد البيانات العلائقية والمتجهية}};
\node[font=\tiny, inner sep=0pt] at (449.3,373.4) {\textarabic{٥}};
\node[font=\tiny, inner sep=0pt] at (467.6,373.4) {\EN{27/10}};
\node[font=\tiny, inner sep=0pt] at (487.1,373.4) {\EN{31/10}};
\node[font=\tiny, inner sep=0pt] at (507.3,373.4) {\EN{BE}};
\node[font=\tiny, inner sep=0pt] at (304.8,363.2) {\textarabic{٢.٢.٣}};
\node[font=\scriptsize, inner sep=0pt] at (378.3,363.2) {\textarabic{الواجهات ولوحة التحكم الإدارية}};
\node[font=\tiny, inner sep=0pt] at (449.3,363.2) {\textarabic{٥}};
\node[font=\tiny, inner sep=0pt] at (467.6,363.2) {\EN{31/10}};
\node[font=\tiny, inner sep=0pt] at (487.1,363.2) {\EN{04/11}};
\node[font=\tiny, inner sep=0pt] at (507.3,363.2) {\EN{FE}};
\node[font=\tiny, inner sep=0pt] at (304.8,352.9) {\textarabic{٢.٢.٤}};
\node[font=\scriptsize, inner sep=0pt] at (378.3,352.9) {\textarabic{مخططات التدفق والتتابع والكيانات}};
\node[font=\tiny, inner sep=0pt] at (449.3,352.9) {\textarabic{٥}};
\node[font=\tiny, inner sep=0pt] at (467.6,352.9) {\EN{03/11}};
\node[font=\tiny, inner sep=0pt] at (487.1,352.9) {\EN{07/11}};
\node[font=\tiny, inner sep=0pt] at (507.3,352.9) {\EN{TM}};
\node[font=\tiny, inner sep=0pt] at (304.8,342.7) {\textarabic{٢.٢.٥}};
\node[font=\scriptsize, inner sep=0pt] at (378.3,342.7) {\textarabic{واجهات برمجة التطبيقات ومواصفاتها}};
\node[font=\tiny, inner sep=0pt] at (449.3,342.7) {\textarabic{٤}};
\node[font=\tiny, inner sep=0pt] at (467.6,342.7) {\EN{06/11}};
\node[font=\tiny, inner sep=0pt] at (487.1,342.7) {\EN{09/11}};
\node[font=\tiny, inner sep=0pt] at (507.3,342.7) {\EN{BE}};
\node[font=\tiny, inner sep=0pt] at (304.8,332.5) {\textarabic{٢.٢.٦}};
\node[font=\scriptsize, inner sep=0pt] at (378.3,332.5) {\textarabic{إطار الأمان والصلاحيات}};
\node[font=\tiny, inner sep=0pt] at (449.3,332.5) {\textarabic{٤}};
\node[font=\tiny, inner sep=0pt] at (467.6,332.5) {\EN{09/11}};
\node[font=\tiny, inner sep=0pt] at (487.1,332.5) {\EN{12/11}};
\node[font=\tiny, inner sep=0pt] at (507.3,332.5) {\EN{BE}};
\node[font=\tiny, inner sep=0pt] at (304.8,322.2) {\textarabic{٢.٢.٧}};
\node[font=\scriptsize, inner sep=0pt] at (378.3,322.2) {\textarabic{مراجعة الفصل الرابع واعتماده}};
\node[font=\tiny, inner sep=0pt] at (449.3,322.2) {\textarabic{٣}};
\node[font=\tiny, inner sep=0pt] at (467.6,322.2) {\EN{12/11}};
\node[font=\tiny, inner sep=0pt] at (487.1,322.2) {\EN{14/11}};
\node[font=\tiny, inner sep=0pt] at (507.3,322.2) {\EN{TM}};
\node[font=\tiny, inner sep=0pt] at (304.8,312.0) {\textarabic{٢.٣}};
\node[font=\tiny\bfseries, inner sep=0pt] at (378.3,312.0) {\textbf{\textarabic{حزمة الفصل الخامس — التنفيذ والتطوير}}};
\node[font=\tiny, inner sep=0pt] at (449.3,312.0) {\textarabic{٣٠}};
\node[font=\tiny, inner sep=0pt] at (467.6,312.0) {\EN{15/11}};
\node[font=\tiny, inner sep=0pt] at (487.1,312.0) {\EN{14/12}};
\node[font=\tiny, inner sep=0pt] at (507.3,312.0) {\textarabic{متعدد}};
\node[font=\tiny, inner sep=0pt] at (304.8,301.8) {\textarabic{٢.٣.١}};
\node[font=\scriptsize, inner sep=0pt] at (378.3,301.8) {\textarabic{د١: بناء وكيل استعلام قواعد البيانات}};
\node[font=\tiny, inner sep=0pt] at (449.3,301.8) {\textarabic{٨}};
\node[font=\tiny, inner sep=0pt] at (467.6,301.8) {\EN{15/11}};
\node[font=\tiny, inner sep=0pt] at (487.1,301.8) {\EN{22/11}};
\node[font=\tiny, inner sep=0pt] at (507.3,301.8) {\EN{AI}};
\node[font=\tiny, inner sep=0pt] at (304.8,291.6) {\textarabic{٢.٣.٢}};
\node[font=\scriptsize, inner sep=0pt] at (378.3,291.6) {\textarabic{د٢: بناء وكيل البحث في المستندات واللوائح}};
\node[font=\tiny, inner sep=0pt] at (449.3,291.6) {\textarabic{٨}};
\node[font=\tiny, inner sep=0pt] at (467.6,291.6) {\EN{23/11}};
\node[font=\tiny, inner sep=0pt] at (487.1,291.6) {\EN{30/11}};
\node[font=\tiny, inner sep=0pt] at (507.3,291.6) {\EN{AI}};
\node[font=\tiny, inner sep=0pt] at (304.8,281.3) {\textarabic{٢.٣.٣}};
\node[font=\scriptsize, inner sep=0pt] at (378.3,281.3) {\textarabic{د٣: دمج الوكلاء وبناء الواجهة الأمامية}};
\node[font=\tiny, inner sep=0pt] at (449.3,281.3) {\textarabic{٨}};
\node[font=\tiny, inner sep=0pt] at (467.6,281.3) {\EN{01/12}};
\node[font=\tiny, inner sep=0pt] at (487.1,281.3) {\EN{08/12}};
\node[font=\tiny, inner sep=0pt] at (507.3,281.3) {\EN{FE}};
\node[font=\tiny, inner sep=0pt] at (304.8,271.1) {\textarabic{٢.٣.٤}};
\node[font=\scriptsize, inner sep=0pt] at (378.3,271.1) {\textarabic{وكلاء التحليل والتقارير — موازٍ}};
\node[font=\tiny, inner sep=0pt] at (449.3,271.1) {\textarabic{١٨}};
\node[font=\tiny, inner sep=0pt] at (467.6,271.1) {\EN{15/11}};
\node[font=\tiny, inner sep=0pt] at (487.1,271.1) {\EN{02/12}};
\node[font=\tiny, inner sep=0pt] at (507.3,271.1) {\EN{BE}};
\node[font=\tiny, inner sep=0pt] at (304.8,260.9) {\textarabic{٢.٣.٥}};
\node[font=\scriptsize, inner sep=0pt] at (378.3,260.9) {\textarabic{التهيئة الذكية للربط والفهرسة — موازٍ}};
\node[font=\tiny, inner sep=0pt] at (449.3,260.9) {\textarabic{١٥}};
\node[font=\tiny, inner sep=0pt] at (467.6,260.9) {\EN{24/11}};
\node[font=\tiny, inner sep=0pt] at (487.1,260.9) {\EN{08/12}};
\node[font=\tiny, inner sep=0pt] at (507.3,260.9) {\EN{BE}};
\node[font=\tiny, inner sep=0pt] at (304.8,250.6) {\textarabic{٢.٣.٦}};
\node[font=\scriptsize, inner sep=0pt] at (378.3,250.6) {\textarabic{سجل التدقيق والإشعارات — موازٍ}};
\node[font=\tiny, inner sep=0pt] at (449.3,250.6) {\textarabic{١٤}};
\node[font=\tiny, inner sep=0pt] at (467.6,250.6) {\EN{01/12}};
\node[font=\tiny, inner sep=0pt] at (487.1,250.6) {\EN{14/12}};
\node[font=\tiny, inner sep=0pt] at (507.3,250.6) {\EN{BE}};
\node[font=\tiny, inner sep=0pt] at (304.8,240.4) {\textarabic{٢.٣.٧}};
\node[font=\scriptsize, inner sep=0pt] at (378.3,240.4) {\textarabic{اختبارات الوحدة والتكامل المستمرة}};
\node[font=\tiny, inner sep=0pt] at (449.3,240.4) {\textarabic{٣٠}};
\node[font=\tiny, inner sep=0pt] at (467.6,240.4) {\EN{15/11}};
\node[font=\tiny, inner sep=0pt] at (487.1,240.4) {\EN{14/12}};
\node[font=\tiny, inner sep=0pt] at (507.3,240.4) {\EN{BE}};
\node[font=\tiny, inner sep=0pt] at (304.8,230.2) {\textarabic{٢.٣.٨}};
\node[font=\scriptsize, inner sep=0pt] at (378.3,230.2) {\textarabic{د٤: الاختبار الشامل وضبط الأداء والنشر}};
\node[font=\tiny, inner sep=0pt] at (449.3,230.2) {\textarabic{٦}};
\node[font=\tiny, inner sep=0pt] at (467.6,230.2) {\EN{09/12}};
\node[font=\tiny, inner sep=0pt] at (487.1,230.2) {\EN{14/12}};
\node[font=\tiny, inner sep=0pt] at (507.3,230.2) {\EN{BE}};
\node[font=\tiny, inner sep=0pt] at (304.8,219.9) {\textarabic{٢.٣.٩}};
\node[font=\scriptsize, inner sep=0pt] at (378.3,219.9) {\textarabic{مراجعة الفصل الخامس واعتماده}};
\node[font=\tiny, inner sep=0pt] at (449.3,219.9) {\textarabic{٣}};
\node[font=\tiny, inner sep=0pt] at (467.6,219.9) {\EN{12/12}};
\node[font=\tiny, inner sep=0pt] at (487.1,219.9) {\EN{14/12}};
\node[font=\tiny, inner sep=0pt] at (507.3,219.9) {\EN{TM}};
\node[font=\tiny, inner sep=0pt] at (304.8,209.7) {\textarabic{٢.٤}};
\node[font=\tiny\bfseries, inner sep=0pt] at (378.3,209.7) {\textbf{\textarabic{حزمة الفصل السادس — الاختبار والنتائج}}};
\node[font=\tiny, inner sep=0pt] at (449.3,209.7) {\textarabic{٢٩}};
\node[font=\tiny, inner sep=0pt] at (467.6,209.7) {\EN{15/12}};
\node[font=\tiny, inner sep=0pt] at (487.1,209.7) {\EN{12/01}};
\node[font=\tiny, inner sep=0pt] at (507.3,209.7) {\textarabic{متعدد}};
\node[font=\tiny, inner sep=0pt] at (304.8,199.5) {\textarabic{٢.٤.١}};
\node[font=\scriptsize, inner sep=0pt] at (378.3,199.5) {\textarabic{اختبار دقة التحويل وزمن الاستجابة}};
\node[font=\tiny, inner sep=0pt] at (449.3,199.5) {\textarabic{٧}};
\node[font=\tiny, inner sep=0pt] at (467.6,199.5) {\EN{15/12}};
\node[font=\tiny, inner sep=0pt] at (487.1,199.5) {\EN{21/12}};
\node[font=\tiny, inner sep=0pt] at (507.3,199.5) {\EN{AI}};
\node[font=\tiny, inner sep=0pt] at (304.8,189.3) {\textarabic{٢.٤.٢}};
\node[font=\scriptsize, inner sep=0pt] at (378.3,189.3) {\textarabic{جودة الاسترجاع والإسناد بالمقاييس}};
\node[font=\tiny, inner sep=0pt] at (449.3,189.3) {\textarabic{٧}};
\node[font=\tiny, inner sep=0pt] at (467.6,189.3) {\EN{22/12}};
\node[font=\tiny, inner sep=0pt] at (487.1,189.3) {\EN{28/12}};
\node[font=\tiny, inner sep=0pt] at (507.3,189.3) {\EN{AI}};
\node[font=\tiny, inner sep=0pt] at (304.8,179.0) {\textarabic{٢.٤.٣}};
\node[font=\scriptsize, inner sep=0pt] at (378.3,179.0) {\textarabic{اختبار الحمل والتحمل}};
\node[font=\tiny, inner sep=0pt] at (449.3,179.0) {\textarabic{٦}};
\node[font=\tiny, inner sep=0pt] at (467.6,179.0) {\EN{29/12}};
\node[font=\tiny, inner sep=0pt] at (487.1,179.0) {\EN{03/01}};
\node[font=\tiny, inner sep=0pt] at (507.3,179.0) {\EN{BE}};
\node[font=\tiny, inner sep=0pt] at (304.8,168.8) {\textarabic{٢.٤.٤}};
\node[font=\scriptsize, inner sep=0pt] at (378.3,168.8) {\textarabic{اختبار الأمان والصلاحيات}};
\node[font=\tiny, inner sep=0pt] at (449.3,168.8) {\textarabic{٧}};
\node[font=\tiny, inner sep=0pt] at (467.6,168.8) {\EN{30/12}};
\node[font=\tiny, inner sep=0pt] at (487.1,168.8) {\EN{05/01}};
\node[font=\tiny, inner sep=0pt] at (507.3,168.8) {\EN{BE}};
\node[font=\tiny, inner sep=0pt] at (304.8,158.6) {\textarabic{٢.٤.٥}};
\node[font=\scriptsize, inner sep=0pt] at (378.3,158.6) {\textarabic{اختبار قبول المستخدمين}};
\node[font=\tiny, inner sep=0pt] at (449.3,158.6) {\textarabic{٥}};
\node[font=\tiny, inner sep=0pt] at (467.6,158.6) {\EN{04/01}};
\node[font=\tiny, inner sep=0pt] at (487.1,158.6) {\EN{08/01}};
\node[font=\tiny, inner sep=0pt] at (507.3,158.6) {\EN{FE}};
\node[font=\tiny, inner sep=0pt] at (304.8,148.3) {\textarabic{٢.٤.٦}};
\node[font=\scriptsize, inner sep=0pt] at (378.3,148.3) {\textarabic{تجميع النتائج وصياغة الفصل السادس}};
\node[font=\tiny, inner sep=0pt] at (449.3,148.3) {\textarabic{٦}};
\node[font=\tiny, inner sep=0pt] at (467.6,148.3) {\EN{05/01}};
\node[font=\tiny, inner sep=0pt] at (487.1,148.3) {\EN{10/01}};
\node[font=\tiny, inner sep=0pt] at (507.3,148.3) {\EN{TM}};
\node[font=\tiny, inner sep=0pt] at (304.8,138.1) {\textarabic{٢.٤.٧}};
\node[font=\scriptsize, inner sep=0pt] at (378.3,138.1) {\textarabic{مراجعة الفصل السادس واعتماده}};
\node[font=\tiny, inner sep=0pt] at (449.3,138.1) {\textarabic{٣}};
\node[font=\tiny, inner sep=0pt] at (467.6,138.1) {\EN{10/01}};
\node[font=\tiny, inner sep=0pt] at (487.1,138.1) {\EN{12/01}};
\node[font=\tiny, inner sep=0pt] at (507.3,138.1) {\EN{TM}};
\node[font=\tiny, inner sep=0pt] at (304.8,127.9) {\textarabic{٢.٥}};
\node[font=\tiny\bfseries, inner sep=0pt] at (378.3,127.9) {\textbf{\textarabic{حزمة الفصل السابع — الخاتمة}}};
\node[font=\tiny, inner sep=0pt] at (449.3,127.9) {\textarabic{١٢}};
\node[font=\tiny, inner sep=0pt] at (467.6,127.9) {\EN{01/01}};
\node[font=\tiny, inner sep=0pt] at (487.1,127.9) {\EN{12/01}};
\node[font=\tiny, inner sep=0pt] at (507.3,127.9) {\textarabic{الفريق}};
\node[font=\tiny, inner sep=0pt] at (304.8,117.6) {\textarabic{٢.٥.١}};
\node[font=\scriptsize, inner sep=0pt] at (378.3,117.6) {\textarabic{الخاتمة وملخص تحقيق الأهداف}};
\node[font=\tiny, inner sep=0pt] at (449.3,117.6) {\textarabic{٤}};
\node[font=\tiny, inner sep=0pt] at (467.6,117.6) {\EN{01/01}};
\node[font=\tiny, inner sep=0pt] at (487.1,117.6) {\EN{04/01}};
\node[font=\tiny, inner sep=0pt] at (507.3,117.6) {\EN{AI}};
\node[font=\tiny, inner sep=0pt] at (304.8,107.4) {\textarabic{٢.٥.٢}};
\node[font=\scriptsize, inner sep=0pt] at (378.3,107.4) {\textarabic{المساهمات والقيود التنفيذية}};
\node[font=\tiny, inner sep=0pt] at (449.3,107.4) {\textarabic{٣}};
\node[font=\tiny, inner sep=0pt] at (467.6,107.4) {\EN{04/01}};
\node[font=\tiny, inner sep=0pt] at (487.1,107.4) {\EN{06/01}};
\node[font=\tiny, inner sep=0pt] at (507.3,107.4) {\EN{TM}};
\node[font=\tiny, inner sep=0pt] at (304.8,97.2) {\textarabic{٢.٥.٣}};
\node[font=\scriptsize, inner sep=0pt] at (378.3,97.2) {\textarabic{التوصيات والأعمال المستقبلية}};
\node[font=\tiny, inner sep=0pt] at (449.3,97.2) {\textarabic{٣}};
\node[font=\tiny, inner sep=0pt] at (467.6,97.2) {\EN{06/01}};
\node[font=\tiny, inner sep=0pt] at (487.1,97.2) {\EN{08/01}};
\node[font=\tiny, inner sep=0pt] at (507.3,97.2) {\EN{TM}};
\node[font=\tiny, inner sep=0pt] at (304.8,87.0) {\textarabic{٢.٥.٤}};
\node[font=\scriptsize, inner sep=0pt] at (378.3,87.0) {\textarabic{مراجعة الفصل السابع واعتماده}};
\node[font=\tiny, inner sep=0pt] at (449.3,87.0) {\textarabic{٣}};
\node[font=\tiny, inner sep=0pt] at (467.6,87.0) {\EN{10/01}};
\node[font=\tiny, inner sep=0pt] at (487.1,87.0) {\EN{12/01}};
\node[font=\tiny, inner sep=0pt] at (507.3,87.0) {\EN{TM}};
\node[font=\tiny, inner sep=0pt] at (304.8,76.7) {\textarabic{٢.٦}};
\node[font=\tiny\bfseries, inner sep=0pt] at (378.3,76.7) {\textbf{\textarabic{حزمة الإصدار النهائي والتسليم}}};
\node[font=\tiny, inner sep=0pt] at (449.3,76.7) {\textarabic{٢٣}};
\node[font=\tiny, inner sep=0pt] at (467.6,76.7) {\EN{01/01}};
\node[font=\tiny, inner sep=0pt] at (487.1,76.7) {\EN{23/01}};
\node[font=\tiny, inner sep=0pt] at (507.3,76.7) {\textarabic{الفريق}};
\node[font=\tiny, inner sep=0pt] at (304.8,66.5) {\textarabic{٢.٦.١}};
\node[font=\scriptsize, inner sep=0pt] at (378.3,66.5) {\textarabic{التدقيق الشامل للوثيقة النهائية}};
\node[font=\tiny, inner sep=0pt] at (449.3,66.5) {\textarabic{٥}};
\node[font=\tiny, inner sep=0pt] at (467.6,66.5) {\EN{01/01}};
\node[font=\tiny, inner sep=0pt] at (487.1,66.5) {\EN{05/01}};
\node[font=\tiny, inner sep=0pt] at (507.3,66.5) {\EN{TM}};
\node[font=\tiny, inner sep=0pt] at (304.8,56.3) {\textarabic{٢.٦.٢}};
\node[font=\scriptsize, inner sep=0pt] at (378.3,56.3) {\textarabic{إعداد البوستر العلمي}};
\node[font=\tiny, inner sep=0pt] at (449.3,56.3) {\textarabic{٤}};
\node[font=\tiny, inner sep=0pt] at (467.6,56.3) {\EN{04/01}};
\node[font=\tiny, inner sep=0pt] at (487.1,56.3) {\EN{07/01}};
\node[font=\tiny, inner sep=0pt] at (507.3,56.3) {\EN{FE}};
\node[font=\tiny, inner sep=0pt] at (304.8,46.0) {\textarabic{٢.٦.٣}};
\node[font=\scriptsize, inner sep=0pt] at (378.3,46.0) {\textarabic{إنتاج الفيديو التسويقي}};
\node[font=\tiny, inner sep=0pt] at (449.3,46.0) {\textarabic{٤}};
\node[font=\tiny, inner sep=0pt] at (467.6,46.0) {\EN{06/01}};
\node[font=\tiny, inner sep=0pt] at (487.1,46.0) {\EN{09/01}};
\node[font=\tiny, inner sep=0pt] at (507.3,46.0) {\EN{FE}};
\node[font=\tiny, inner sep=0pt] at (304.8,35.8) {\textarabic{٢.٦.٤}};
\node[font=\scriptsize, inner sep=0pt] at (378.3,35.8) {\textarabic{التسليم النهائي الشامل للكلية}};
\node[font=\tiny, inner sep=0pt] at (449.3,35.8) {\textarabic{٢}};
\node[font=\tiny, inner sep=0pt] at (467.6,35.8) {\EN{08/01}};
\node[font=\tiny, inner sep=0pt] at (487.1,35.8) {\EN{09/01}};
\node[font=\tiny, inner sep=0pt] at (507.3,35.8) {\EN{TM}};
\node[font=\tiny, inner sep=0pt] at (304.8,25.6) {\textarabic{٢.٦.٥}};
\node[font=\scriptsize, inner sep=0pt] at (378.3,25.6) {\textarabic{العرض التقديمي النهائي والمحاكاة}};
\node[font=\tiny, inner sep=0pt] at (449.3,25.6) {\textarabic{٦}};
\node[font=\tiny, inner sep=0pt] at (467.6,25.6) {\EN{10/01}};
\node[font=\tiny, inner sep=0pt] at (487.1,25.6) {\EN{15/01}};
\node[font=\tiny, inner sep=0pt] at (507.3,25.6) {\EN{FE}};
\node[font=\tiny, inner sep=0pt] at (304.8,15.3) {\textarabic{٢.٦.٦}};
\node[font=\scriptsize, inner sep=0pt] at (378.3,15.3) {\textarabic{المناقشة النهائية — بوابة الجودة الثانية}};
\node[font=\tiny, inner sep=0pt] at (449.3,15.3) {\textarabic{٨}};
\node[font=\tiny, inner sep=0pt] at (467.6,15.3) {\EN{16/01}};
\node[font=\tiny, inner sep=0pt] at (487.1,15.3) {\EN{23/01}};
\node[font=\tiny, inner sep=0pt] at (507.3,15.3) {\EN{TM}};
\node[font=\tiny, inner sep=0pt] at (304.8,5.1) {\textarabic{٢.٦.٧}};
\node[font=\scriptsize, inner sep=0pt] at (378.3,5.1) {\textarabic{الأرشفة الختامية للمشروع}};
\node[font=\tiny, inner sep=0pt] at (449.3,5.1) {\textarabic{٤}};
\node[font=\tiny, inner sep=0pt] at (467.6,5.1) {\EN{20/01}};
\node[font=\tiny, inner sep=0pt] at (487.1,5.1) {\EN{23/01}};
\node[font=\tiny, inner sep=0pt] at (507.3,5.1) {\EN{TM}};
\end{tikzpicture}%
\end{LTR}%
\end{lrbox}%
\noindent\makebox[\textwidth][c]{\usebox\gCBB}%
\begin{center}
\scriptsize
\renewcommand{\arraystretch}{1.0}
\begin{tabular}{c@{\hspace{3pt}}l@{\hspace{7pt}}c@{\hspace{3pt}}l@{\hspace{7pt}}c@{\hspace{3pt}}l}
  \tikz\fill[ownerAI!55,draw=ownerAI] (0,0) rectangle (0.28,0.14); & الذكاء الاصطناعي (\EN{AI}) &
  \tikz\fill[ownerBE!55,draw=ownerBE] (0,0) rectangle (0.28,0.14); & الخلفية وقواعد البيانات (\EN{BE}) &
  \tikz\fill[ownerFE!45,draw=ownerFE] (0,0) rectangle (0.28,0.14); & الواجهة وتجربة المستخدم (\EN{FE}) \\[3pt]
  \tikz\fill[ownerTM!50,draw=ownerTM] (0,0) rectangle (0.28,0.14); & الفريق كاملاً (\EN{TM}) &
  \tikz\fill[kaudark!90,draw=kaudark!90] (0,0.05) rectangle (0.4,0.11); & شريط الحزمة (الملخص) &
  \tikz{\fill[kauaccent,draw=kauaccent!50!black] (0.13,0.26) -- (0.26,0.13) -- (0.13,0) -- (0,0.13) -- cycle;} & معلم رسمي بتاريخه الرسمي \\[3pt]
  \tikz\fill[kauaccent!35,draw=kauaccent!70!black] (0,0) rectangle (0.28,0.14); & إعداد العرض والمحاكاة &
  \tikz\fill[kauaccent!75,draw=kauaccent!80!black] (0,0) rectangle (0.28,0.14); & أسبوع المناقشة النهائية &
  \tikz{\draw[-stealth,thick,kauprimary] (0,0.07) -- (0.4,0.07);} & سهم تبعية (المسار الحرج) \\[3pt]
  \tikz\fill[gray!40] (0,0.02) rectangle (0.3,0.13); & ظل عطلة نهاية الأسبوع &
  \tikz{\draw[gray!45] (0.02,0) -- (0.02,0.16);} & بداية الأسبوع (الأحد) &
  \tikz\fill[kauprimary!22,draw=kauprimary!55] (0,0) rectangle (0.3,0.14); & ترويسة لوحة البيانات \\
\end{tabular}
\end{center}
  \caption{مخطط جانت التنفيذي لمشروع التخرج 2 — المرحلتان الثانية والثالثة (أكتوبر 2026 -- يناير 2027): لوحة موحَّدة لـ39 مهمة مع الدورات التطويرية الأربع والمعالم \EN{M8}--\EN{M13}.}
  \label{fig:ganttP2}
\end{figure}


ويُقرأ لون شريط المهمة وفق مفتاح الرموز أسفل كل مخطط: الأزرق المخضرّ
للذكاء الاصطناعي، والأزرق للواجهة الخلفية وقواعد البيانات، والبنفسجي
للواجهة الأمامية وتجربة المستخدم، والرمادي للفريق كاملاً؛ بينما تحتفظ أشرطة إعداد العرض
التقديمي والمحاكاة الداخلية بلونها الأحمر الفاتح المميز لنوافذ ما قبل
المناقشة، ويُبرز الشريط الأحمر الغامق أسبوع المناقشة النهائية بوصفه بوابة
الجودة الختامية لكل مرحلة. وتُثبَّت المعالم الرسمية بمعيناتها عند
تواريخها في مواضعها من صفوف مهامها، وتجسد الأسهم
الزرقاء تسلسل المسار الحرج الذي لا يحتمل التأخير.

ويعكس المخطط الثاني بنية المرحلة الثانية الأجايلية بوضوح؛ إذ تظهر الدورات
التطويرية الأربع (\EN{Sprint 1}--\EN{Sprint 4}) متتابعةً في مقاطع زمنية
مستقلة تنتهي كل منها بمراجعة اعتماد المخرج مع المشرف، بينما تمتد الأشرطة
الموازية (وكلاء التحليل والتقارير، والتهيئة الذكية، وسجل التدقيق،
والاختبارات المستمرة) على امتداد نافذة التطوير كاملة لتجسيد مبدأ العمل
المتوازي الذي تتيحه المنهجية الرشيقة، وهو ما يمنح الفريق قدرة أكبر على
امتصاص أي انحراف جزئي دون المساس بمواعيد البوابات الرسمية الثابتة.

%------------------------------------------------------------------------------
\section{متطلبات الأدوات (\EN{Required Tools})}
\label{sec:tools}

يستند تطوير مشروع بيان إلى مجموعة من الأدوات المادية والبرمجية التي تضمن
جودة المنتج النهائي وكفاءة عملية التطوير. وفيما يلي عرض تفصيلي لهذه الأدوات.

\subsection{الأدوات المادية (\EN{Hardware Tools})}
\label{subsec:hardware}

يوضح الجدول~\ref{tab:hardware} المواصفات المادية المطلوبة لكل عضو في فريق
التطوير، إلى جانب الأجهزة اللازمة لتشغيل بيئة الخوادم.

\begin{table}[H]
  \centering
  \caption{الأدوات المادية المطلوبة للمشروع.}
  \label{tab:hardware}
  \small
  \renewcommand{\arraystretch}{1.4}
  \begin{tabularx}{\textwidth}{|>{\raggedright\arraybackslash}p{2.6cm}|>{\centering\arraybackslash}p{1.2cm}|>{\raggedright\arraybackslash}p{4.4cm}|>{\raggedright\arraybackslash}X|}
    \hline
    \textbf{الأداة} & \textbf{الكمية} & \textbf{الاستخدام} & \textbf{المواصفات} \\
    \hline
    اتصال بالإنترنت & --- & للبحث، وتصحيح الأخطاء، وتحميل النماذج والأدوات، والتواصل مع الفريق & سرعة مناسبة للبحث والعمل الجماعي وتحميل ملفات بحجم عدة جيجابايتات (مثل نماذج \EN{Ollama}). \\
    \hline
    لابتوب & 5 & لكل عضو جهاز للعمل على تطوير المشروع &
      \textbf{المعالج:} \EN{Core i7} أو ما يعادله (لتشغيل بيئة التطوير).\newline
      \textbf{الرام} (\EN{RAM}): 16 \EN{GB} كحد أدنى، ويفضل 32 \EN{GB} لتشغيل النماذج المحلية.\newline
      \textbf{مساحة التخزين:} 512 \EN{GB SSD} أو أكثر (لأنظمة التشغيل، قواعد البيانات، ونماذج الذكاء الاصطناعي). \\
    \hline
    خادم محلي أو حاسوب إضافي & 1 & تشغيل خدمات الخلفية الأساسية مثل قواعد البيانات (\EN{PostgreSQL}، \EN{Qdrant}، \EN{Redis}) ونماذج \EN{Ollama} محلياً &
      \textbf{المعالج:} \EN{Core i7} أو ما يعادله مع دعم لتسريع \EN{CUDA} (بطاقة رسوميات \EN{NVIDIA}).\newline
      \textbf{الرام} (\EN{RAM}): 32 \EN{GB} أو أكثر.\newline
      \textbf{مساحة التخزين:} 1 \EN{TB SSD} أو أكثر.\newline
      \textbf{نظام التشغيل:} \EN{Linux} (\EN{Ubuntu}) أو \EN{Windows} مع \EN{WSL2}. \\
    \hline
    ذاكرة فلاش (\EN{Flash memory}) & 2 & لعمل نسخ احتياطية للمشروع، وملفات التطوير، والبيانات الأولية & سعة 64 \EN{GB} أو أكثر، ويفضل أن تدعم \EN{USB 3.0} للسرعة. \\
    \hline
  \end{tabularx}
\end{table}

\subsection{الأدوات البرمجية (\EN{Software Tools})}
\label{subsec:software}

يستعرض الجدول~\ref{tab:software} الأدوات البرمجية الرئيسية المعتمدة في تطوير
المشروع، مع ذكر الاستخدام المخصص لكل أداة.

\begin{longtable}{|>{\raggedright\arraybackslash}p{5.5cm}|>{\raggedright\arraybackslash}p{8cm}|}
\caption{الأدوات البرمجية المعتمدة في تطوير مشروع بيان.}
\label{tab:software}\\
\hline
\textbf{البرنامج (\EN{Software})} & \textbf{الاستخدام} \\
\hline
\endfirsthead
\multicolumn{2}{c}{جدول~\thetable{} — تابع}\\
\hline
\textbf{البرنامج (\EN{Software})} & \textbf{الاستخدام} \\
\hline
\endhead

\EN{Windows 11} & نظام تشغيل لتشغيل النظام على حواسيب التطوير، مع دعم \EN{WSL2} لتشغيل بيئات \EN{Linux} على \EN{Windows}. \\
\hline
\EN{LaTeX} (\EN{TeX Live} + \EN{latexmk}) & لكتابة وتوثيق تقرير المشروع الأكاديمي. \\
\hline
\EN{Microsoft Project 2021} & لإعداد الجدول الزمني للمشروع ومخططات جانت. \\
\hline
\EN{Microsoft PowerPoint 2021} & لإعداد العرض التقديمي للمشروع. \\
\hline
\EN{Draw.io} & لإنشاء المخططات والنماذج الخاصة بالمشروع. \\
\hline
\EN{Figma} & لتصميم واجهات النظام (\EN{UI/UX}). \\
\hline
\EN{Visual Studio Code} & بيئة التطوير المتكاملة لكتابة كود المشروع (\EN{Python}، \EN{JavaScript}، \EN{TypeScript}). \\
\hline
\EN{Python 3.12+} & لغة البرمجة الأساسية لتطوير الواجهة الخلفية (\EN{FastAPI}) ووكلاء \EN{LangGraph}. \\
\hline
\EN{Node.js 18+} & بيئة تشغيل (\EN{Runtime}) لتطوير الواجهة الأمامية (\EN{Next.js}) وإدارة الحزم عبر \EN{npm}. \\
\hline
\EN{FastAPI} & إطار عمل لبناء الواجهة الخلفية (\EN{Backend API}) وخدمات الوكلاء. \\
\hline
\EN{Next.js} & إطار عمل لبناء الواجهة الأمامية (\EN{Frontend}) مع دعم \EN{RTL} للغة العربية. \\
\hline
\EN{LangChain} / \EN{LangGraph} & مكتبات لبناء وتنسيق الوكلاء المتعددين (\EN{Multi-Agent System}). \\
\hline
\EN{PostgreSQL} مع \EN{pgvector} & قاعدة بيانات علائقية لتخزين البيانات المؤسسية والمتجهات (لـ \EN{RAG}). \\
\hline
\EN{Qdrant} & قاعدة بيانات متجهة (\EN{Vector Database}) لتخزين واسترجاع تضمينات المستندات. \\
\hline
\EN{Redis} & تخزين مؤقت (\EN{Cache}) وطابور رسائل لتنسيق الوكلاء وإدارة الجلسات. \\
\hline
\EN{Ollama} & أداة لتشغيل النماذج اللغوية محلياً (مثل \EN{Qwen}، \EN{Gemma}، \EN{Llama}). \\
\hline
\EN{Docker Desktop} & منصة للحاويات لتشغيل جميع الخدمات (\EN{PostgreSQL}، \EN{Qdrant}، \EN{Redis}، \EN{Ollama}) بشكل موحد. \\
\hline
\EN{Git} / \EN{GitHub} & نظام التحكم في الإصدارات وإدارة المستودع البرمجي للمشروع. \\
\hline
\EN{Postman / Insomnia} & اختبار واجهات برمجة التطبيقات (\EN{APIs}) للواجهة الخلفية. \\
\hline

\end{longtable}

%------------------------------------------------------------------------------
\section{هيكل المشروع (\EN{Project Structure})}
\label{sec:structure}

يتكون هذا المشروع من سبعة (7) فصول رئيسية، يُكمل كلٌّ منها ما سبقه:

\begin{itemize}
  \item \textbf{الفصل الأول — المقدمة:} الخلفية، المشكلة، الأهداف، الأهمية، المساهمات،
        النطاق، المنهجية، الجدول الزمني، والأدوات.
  \item \textbf{الفصل الثاني — الخلفية النظرية ومراجعة الأدبيات:} الأسس
        النظرية للتقنيات الأربع، الأنظمة السابقة، المقارنة، والفجوة البحثية.
  \item \textbf{الفصل الثالث — المنهجية وتحليل المشروع:} منهجية التطوير، دراسة
        الجدوى، تحليل المتطلبات، حالات
        الاستخدام، وتحليل المخاطر.
  \item \textbf{الفصل الرابع — تصميم النظام:} المعمارية العامة، تصميم قواعد
        البيانات، تصميم واجهات المستخدم، المخططات، وتصميم \EN{APIs}.
  \item \textbf{الفصل الخامس — تنفيذ واختبار النظام:} بيئة التطوير، التنفيذ،
        النشر، واختبارات (الوحدة، التكامل، الأداء، وقبول المستخدم).
  \item \textbf{الفصل السادس — النتائج والمناقشة:} النتائج الكمية والنوعية،
        واجهات النظام النهائية، المناقشة، التحديات، والدروس المستفادة.
  \item \textbf{الفصل السابع — الخاتمة والأعمال المستقبلية:} الخاتمة، تحقيق
        الأهداف، المساهمات، القيود، التوصيات، والأعمال المستقبلية.
\end{itemize}
